% Options for packages loaded elsewhere
% Options for packages loaded elsewhere
\PassOptionsToPackage{unicode}{hyperref}
\PassOptionsToPackage{hyphens}{url}
%
\documentclass[
  ignorenonframetext,
]{beamer}
\newif\ifbibliography
\usepackage{pgfpages}
\setbeamertemplate{caption}[numbered]
\setbeamertemplate{caption label separator}{: }
\setbeamercolor{caption name}{fg=normal text.fg}
\beamertemplatenavigationsymbolsempty
% remove section numbering
\setbeamertemplate{part page}{
  \centering
  \begin{beamercolorbox}[sep=16pt,center]{part title}
    \usebeamerfont{part title}\insertpart\par
  \end{beamercolorbox}
}
\setbeamertemplate{section page}{
  \centering
  \begin{beamercolorbox}[sep=12pt,center]{section title}
    \usebeamerfont{section title}\insertsection\par
  \end{beamercolorbox}
}
\setbeamertemplate{subsection page}{
  \centering
  \begin{beamercolorbox}[sep=8pt,center]{subsection title}
    \usebeamerfont{subsection title}\insertsubsection\par
  \end{beamercolorbox}
}
% Prevent slide breaks in the middle of a paragraph
\widowpenalties 1 10000
\raggedbottom
\AtBeginPart{
  \frame{\partpage}
}
\AtBeginSection{
  \ifbibliography
  \else
    \frame{\sectionpage}
  \fi
}
\AtBeginSubsection{
  \frame{\subsectionpage}
}
\usepackage{iftex}
\ifPDFTeX
  \usepackage[T1]{fontenc}
  \usepackage[utf8]{inputenc}
  \usepackage{textcomp} % provide euro and other symbols
\else % if luatex or xetex
  \usepackage{unicode-math} % this also loads fontspec
  \defaultfontfeatures{Scale=MatchLowercase}
  \defaultfontfeatures[\rmfamily]{Ligatures=TeX,Scale=1}
\fi
\usepackage{lmodern}

\ifPDFTeX\else
  % xetex/luatex font selection
\fi
% Use upquote if available, for straight quotes in verbatim environments
\IfFileExists{upquote.sty}{\usepackage{upquote}}{}
\IfFileExists{microtype.sty}{% use microtype if available
  \usepackage[]{microtype}
  \UseMicrotypeSet[protrusion]{basicmath} % disable protrusion for tt fonts
}{}
\makeatletter
\@ifundefined{KOMAClassName}{% if non-KOMA class
  \IfFileExists{parskip.sty}{%
    \usepackage{parskip}
  }{% else
    \setlength{\parindent}{0pt}
    \setlength{\parskip}{6pt plus 2pt minus 1pt}}
}{% if KOMA class
  \KOMAoptions{parskip=half}}
\makeatother


\usepackage{longtable,booktabs,array}
\usepackage{calc} % for calculating minipage widths
\usepackage{caption}
% Make caption package work with longtable
\makeatletter
\def\fnum@table{\tablename~\thetable}
\makeatother
\usepackage{graphicx}
\makeatletter
\newsavebox\pandoc@box
\newcommand*\pandocbounded[1]{% scales image to fit in text height/width
  \sbox\pandoc@box{#1}%
  \Gscale@div\@tempa{\textheight}{\dimexpr\ht\pandoc@box+\dp\pandoc@box\relax}%
  \Gscale@div\@tempb{\linewidth}{\wd\pandoc@box}%
  \ifdim\@tempb\p@<\@tempa\p@\let\@tempa\@tempb\fi% select the smaller of both
  \ifdim\@tempa\p@<\p@\scalebox{\@tempa}{\usebox\pandoc@box}%
  \else\usebox{\pandoc@box}%
  \fi%
}
% Set default figure placement to htbp
\def\fps@figure{htbp}
\makeatother





\setlength{\emergencystretch}{3em} % prevent overfull lines

\providecommand{\tightlist}{%
  \setlength{\itemsep}{0pt}\setlength{\parskip}{0pt}}



 


\makeatletter
\@ifpackageloaded{tcolorbox}{}{\usepackage[skins,breakable]{tcolorbox}}
\@ifpackageloaded{fontawesome5}{}{\usepackage{fontawesome5}}
\definecolor{quarto-callout-color}{HTML}{909090}
\definecolor{quarto-callout-note-color}{HTML}{0758E5}
\definecolor{quarto-callout-important-color}{HTML}{CC1914}
\definecolor{quarto-callout-warning-color}{HTML}{EB9113}
\definecolor{quarto-callout-tip-color}{HTML}{00A047}
\definecolor{quarto-callout-caution-color}{HTML}{FC5300}
\definecolor{quarto-callout-color-frame}{HTML}{acacac}
\definecolor{quarto-callout-note-color-frame}{HTML}{4582ec}
\definecolor{quarto-callout-important-color-frame}{HTML}{d9534f}
\definecolor{quarto-callout-warning-color-frame}{HTML}{f0ad4e}
\definecolor{quarto-callout-tip-color-frame}{HTML}{02b875}
\definecolor{quarto-callout-caution-color-frame}{HTML}{fd7e14}
\makeatother
\makeatletter
\@ifpackageloaded{caption}{}{\usepackage{caption}}
\AtBeginDocument{%
\ifdefined\contentsname
  \renewcommand*\contentsname{Table of contents}
\else
  \newcommand\contentsname{Table of contents}
\fi
\ifdefined\listfigurename
  \renewcommand*\listfigurename{List of Figures}
\else
  \newcommand\listfigurename{List of Figures}
\fi
\ifdefined\listtablename
  \renewcommand*\listtablename{List of Tables}
\else
  \newcommand\listtablename{List of Tables}
\fi
\ifdefined\figurename
  \renewcommand*\figurename{Figure}
\else
  \newcommand\figurename{Figure}
\fi
\ifdefined\tablename
  \renewcommand*\tablename{Table}
\else
  \newcommand\tablename{Table}
\fi
}
\@ifpackageloaded{float}{}{\usepackage{float}}
\floatstyle{ruled}
\@ifundefined{c@chapter}{\newfloat{codelisting}{h}{lop}}{\newfloat{codelisting}{h}{lop}[chapter]}
\floatname{codelisting}{Listing}
\newcommand*\listoflistings{\listof{codelisting}{List of Listings}}
\makeatother
\makeatletter
\makeatother
\makeatletter
\@ifpackageloaded{caption}{}{\usepackage{caption}}
\@ifpackageloaded{subcaption}{}{\usepackage{subcaption}}
\makeatother

\usepackage{bookmark}
\IfFileExists{xurl.sty}{\usepackage{xurl}}{} % add URL line breaks if available
\urlstyle{same}
\hypersetup{
  pdftitle={Tema 8: Flujos en redes},
  pdfauthor={Víctor Aceña Gil - Antonio Alonso Ayuso},
  hidelinks,
  pdfcreator={LaTeX via pandoc}}


\title{Tema 8: Flujos en redes}
\subtitle{Optimización y Análisis de Redes}
\author{Víctor Aceña Gil - Antonio Alonso Ayuso}
\date{2026-08-18}

\begin{document}
\frame{\titlepage}


\begin{frame}{Objetivos de aprendizaje}
\phantomsection\label{objetivos-de-aprendizaje}
\begin{enumerate}
\tightlist
\item
  \textbf{Formular} algebraicamente el concepto de flujo y circulación
  en digrafos mediante las ecuaciones de conservación y balance en los
  nodos.
\item
  \textbf{Caracterizar} la estructura vectorial de los flujos,
  calculando las bases fundamentales de ciclos y el número ciclomático
  \(\nu(G)\).
\item
  \textbf{Demostrar} y \textbf{aplicar} el Lema de Coloración de Arcos
  de Minty mediante su algoritmo constructivo de etiquetado de vértices.
\item
  \textbf{Demostrar} y \textbf{analizar} el Teorema de Flujo Máximo -
  Corte Mínimo de Ford-Fulkerson y la relación de dualidad en
  optimización lineal.
\item
  \textbf{Ejecutar} y \textbf{caracterizar} el algoritmo de
  Ford-Fulkerson en redes residuales mediante el método de etiquetado
  directo y la cota de Edmonds-Karp.
\item
  \textbf{Formular} y \textbf{resolver} los problemas de flujo a coste
  mínimo y el problema general del transbordo.
\end{enumerate}
\end{frame}

\section{Estructura algebraica de flujos, ciclos y
cociclos}\label{estructura-algebraica-de-flujos-ciclos-y-cociclos}

\begin{frame}{Introducción y contexto físico}
\phantomsection\label{introducciuxf3n-y-contexto-fuxedsico}
\begin{itemize}
\tightlist
\item
  \textbf{Representación de redes de flujo}:

  \begin{itemize}
  \tightlist
  \item
    Digrafo conexo \(G = (V, U)\) con \(n = |V|\) vértices y \(m = |U|\)
    arcos dirigidos.
  \item
    Modela la circulación de bienes (tráfico vehicular, fluidos,
    información, electricidad, capitales).
  \end{itemize}
\item
  \textbf{Leyes físicas y topológicas}:

  \begin{itemize}
  \tightlist
  \item
    Los flujos responden a leyes de conservación local en los vértices
    (análogas a las leyes de Kirchhoff en circuitos o continuidad en
    fluidos).
  \end{itemize}
\item
  \textbf{Estructura matemática del problema}:

  \begin{itemize}
  \tightlist
  \item
    Álgebra lineal aplicada sobre la topología del digrafo.
  \item
    Dualidad estricta entre circulaciones cerradas (ciclos) y barreras
    de corte (cociclos).
  \end{itemize}
\end{itemize}
\end{frame}

\begin{frame}{Definición: Flujo en un digrafo y conservación}
\phantomsection\label{definiciuxf3n-flujo-en-un-digrafo-y-conservaciuxf3n}
\begin{tcolorbox}[enhanced jigsaw, coltitle=black, toprule=.15mm, opacitybacktitle=0.6, toptitle=1mm, leftrule=.75mm, breakable, left=2mm, opacityback=0, colback=white, bottomtitle=1mm, titlerule=0mm, title=\textcolor{quarto-callout-note-color}{\faInfo}\hspace{0.5em}{Flujo en un digrafo y conservación}, colbacktitle=quarto-callout-note-color!10!white, rightrule=.15mm, bottomrule=.15mm, arc=.35mm, colframe=quarto-callout-note-color-frame]

Dado un digrafo conexo \(G = (V, U)\) con \(m\) arcos dirigidos, un
\textbf{flujo} es un vector
\(\mathbf{f} = (f_{u_1}, f_{u_2}, \dots, f_{u_m})^T \in \mathbb{R}^m\)
cuyas componentes representan el caudal \(f_u\) que transita por cada
arco \(u \in U\), satisfaciendo la \textbf{ley de conservación de flujo}
en todos los vértices:

\[
\sum_{u \in w^+(i)} f_u = \sum_{u \in w^-(i)} f_u, \quad \forall i \in V
\]

donde \(w^+(i)\) son los arcos salientes del nodo \(i\) y \(w^-(i)\) los
arcos entrantes al nodo \(i\).

\end{tcolorbox}
\end{frame}

\begin{frame}{Nota de notación: Arcos, vectores y componentes}
\phantomsection\label{nota-de-notaciuxf3n-arcos-vectores-y-componentes}
\begin{tcolorbox}[enhanced jigsaw, coltitle=black, toprule=.15mm, opacitybacktitle=0.6, toptitle=1mm, leftrule=.75mm, breakable, left=2mm, opacityback=0, colback=white, bottomtitle=1mm, titlerule=0mm, title=\textcolor{quarto-callout-note-color}{\faInfo}\hspace{0.5em}{Distinción entre arcos, vectores y componentes}, colbacktitle=quarto-callout-note-color!10!white, rightrule=.15mm, bottomrule=.15mm, arc=.35mm, colframe=quarto-callout-note-color-frame]

Para evitar cualquier confusión entre los nombres de los arcos y los
valores numéricos del flujo:

\begin{enumerate}
\tightlist
\item
  \textbf{Identificador del Arco}: Los \(m\) arcos del digrafo se
  etiquetan como \(u_1, u_2, \dots, u_m \in U\) (o por pares ordenados
  \(u = (i, j)\)).
\item
  \textbf{Identificador de Vectores (Superíndices)}: Los distintos
  vectores de flujo o de ciclo se denotan en negrita mediante
  \textbf{superíndices} (\(\mathbf{f}^1, \mathbf{f}^2, \mathbf{f}^3\);
  \(\mathbf{\mu}^1, \mathbf{\mu}^2\)).
\item
  \textbf{Componentes Escalares de Arco (Subíndices)}: El caudal por el
  arco \(u_k\) se denota mediante el \textbf{subíndice} \(f_{u_k}\) (o
  \(f_{ij}\)). En los gráficos, se especifica el arco y su valor
  (\(u_k: f_{u_k}\)).
\end{enumerate}

\end{tcolorbox}
\end{frame}

\begin{frame}{Balances nodales y tipos de nodos}
\phantomsection\label{balances-nodales-y-tipos-de-nodos}
\begin{itemize}
\item
  \textbf{Escalar de balance \(b_i \in \mathbb{R}\) en cada nodo
  \(i \in V\)}:

  \begin{itemize}
  \tightlist
  \item
    \textbf{Nodo fuente (\(b_i > 0\))}: Inyecta \(b_i\) unidades netas
    de flujo a la red.
  \item
    \textbf{Nodo sumidero (\(b_i < 0\))}: Extrae \(|b_i|\) unidades
    netas de flujo de la red.
  \item
    \textbf{Nodo de transbordo (\(b_i = 0\))}: La suma entrante iguala a
    la saliente (nodo de paso).
  \end{itemize}
\item
  \textbf{Condición de equilibrio global}: \[\sum_{i \in V} b_i = 0\]
\item
  \textbf{Concepto de Circulación}:

  \begin{itemize}
  \tightlist
  \item
    Cuando \(b_i = 0\) para todo \(i \in V\), la red no posee entradas
    ni salidas externas.
  \item
    El vector de flujo \(\mathbf{f}\) se denomina estrictamente
    \textbf{circulación}.
  \end{itemize}
\end{itemize}
\end{frame}

\begin{frame}{Ejemplo: Representación gráfica de los tres vectores}
\phantomsection\label{ejemplo-representaciuxf3n-gruxe1fica-de-los-tres-vectores}
\begin{itemize}
\tightlist
\item
  Representación visual de los vectores de flujo \(\mathbf{f}^1\),
  \(\mathbf{f}^2\) y \(\mathbf{f}^3\) evaluados sobre la red de 5 nodos.
\end{itemize}
\end{frame}

\begin{frame}{Ejemplo: Evaluación del vector de flujo f1}
\phantomsection\label{ejemplo-evaluaciuxf3n-del-vector-de-flujo-f1}
Considérese el digrafo de 5 nodos y 7 arcos para evaluar el primer
vector de flujo \(\mathbf{f}^1 \in \mathbb{R}^7\):

\begin{itemize}
\tightlist
\item
  \textbf{Vector
  \(\mathbf{f}^1 = (2.7;\, 2;\, 2;\, 2;\, 2.7;\, -0.7;\, 0)^T\)}:

  \begin{itemize}
  \tightlist
  \item
    Nodo \(a\): Salen \(f_1 = 2.7\), entran
    \(f_5 = 2.7 \implies 2.7 = 2.7\).
  \item
    Nodo \(b\): Salen \(f_2 = 2\), entran
    \(f_1 + f_6 = 2.7 + (-0.7) = 2 \implies 2 = 2\).
  \item
    Nodo \(e\): Salen \(f_5 + f_6 + f_7 = 2.7 - 0.7 + 0 = 2\), entra
    \(f_4 = 2 \implies 2 = 2\).
  \item
    \emph{Conclusión}: Cumple la ley de conservación en todos los nodos
    \(\implies\) \textbf{Flujo válido}.
  \end{itemize}
\end{itemize}
\end{frame}

\begin{frame}{Ejemplo: Evaluación de los vectores f2 y f3}
\phantomsection\label{ejemplo-evaluaciuxf3n-de-los-vectores-f2-y-f3}
Evaluación de los siguientes dos vectores de flujo en la misma red:

\begin{itemize}
\tightlist
\item
  \textbf{Vector
  \(\mathbf{f}^2 = (0;\, 1;\, 1;\, 1;\, 0;\, 1;\, 0)^T\)}:

  \begin{itemize}
  \tightlist
  \item
    Flujo circulante por el ciclo formado por los arcos
    \(u_2, u_3, u_4, u_6\).
  \item
    Satisface la conservación en todos los vértices.
  \end{itemize}
\item
  \textbf{Vector
  \(\mathbf{f}^3 = (1;\, 0;\, -1;\, -1;\, 1;\, -1;\, -1)^T\)}:

  \begin{itemize}
  \tightlist
  \item
    Satisface las ecuaciones de conservación en los 5 nodos.
  \end{itemize}
\end{itemize}
\end{frame}

\begin{frame}{Lema: Flujo nulo en árboles}
\phantomsection\label{lema-flujo-nulo-en-uxe1rboles}
\begin{tcolorbox}[enhanced jigsaw, coltitle=black, toprule=.15mm, opacitybacktitle=0.6, toptitle=1mm, leftrule=.75mm, breakable, left=2mm, opacityback=0, colback=white, bottomtitle=1mm, titlerule=0mm, title=\textcolor{quarto-callout-note-color}{\faInfo}\hspace{0.5em}{Flujo nulo en árboles}, colbacktitle=quarto-callout-note-color!10!white, rightrule=.15mm, bottomrule=.15mm, arc=.35mm, colframe=quarto-callout-note-color-frame]

Un flujo definido sobre un árbol (digrafo conexo acíclico) es
necesariamente el \textbf{flujo nulo} (\(\mathbf{f} = \mathbf{0}\)).

\end{tcolorbox}

\begin{itemize}
\tightlist
\item
  \textbf{Demostración inductiva}:

  \begin{itemize}
  \tightlist
  \item
    Todo árbol de \(n\) vértices posee \(n - 1\) arcos y contiene al
    menos una hoja \(x_1\) de grado 1.
  \item
    Sea \(u\) el único arco incidente en \(x_1\). Por conservación,
    \(f_u = 0\).
  \item
    Retirando inductivamente la hoja \(x_1\), se deduce que todas las
    componentes son nulas: \(f_u = 0, \forall u \in U\).
  \end{itemize}
\item
  \textbf{Conclusión fundamental}: La presencia de \textbf{ciclos
  cerrados es estrictamente necesaria} para sostener circulaciones no
  nulas en una red.
\end{itemize}
\end{frame}

\begin{frame}{Lema: Estructura vectorial del espacio de flujos}
\phantomsection\label{lema-estructura-vectorial-del-espacio-de-flujos}
\begin{tcolorbox}[enhanced jigsaw, coltitle=black, toprule=.15mm, opacitybacktitle=0.6, toptitle=1mm, leftrule=.75mm, breakable, left=2mm, opacityback=0, colback=white, bottomtitle=1mm, titlerule=0mm, title=\textcolor{quarto-callout-note-color}{\faInfo}\hspace{0.5em}{Estructura vectorial del espacio de flujos}, colbacktitle=quarto-callout-note-color!10!white, rightrule=.15mm, bottomrule=.15mm, arc=.35mm, colframe=quarto-callout-note-color-frame]

Cualquier combinación lineal de flujos es un flujo. En consecuencia, el
conjunto de flujos definidos sobre un digrafo \(G = (V, U)\) forma un
espacio vectorial.

\end{tcolorbox}

\begin{itemize}
\tightlist
\item
  \textbf{Clausura algebraica}:

  \begin{itemize}
  \tightlist
  \item
    Si \(\mathbf{f}^1\) y \(\mathbf{f}^2\) son flujos y
    \(\alpha, \beta \in \mathbb{R}\), entonces
    \(\alpha \mathbf{f}^1 + \beta \mathbf{f}^2\) es un flujo factible.
  \item
    La suma y el escalado de circulaciones preservan exactamente la
    conservación nodal.
  \end{itemize}
\item
  \textbf{Objetivo de caracterización}:

  \begin{itemize}
  \tightlist
  \item
    Identificar la dimensión de este espacio vectorial y construir una
    \textbf{base de flujos}.
  \end{itemize}
\end{itemize}
\end{frame}

\begin{frame}{Ejemplo: Representación del flujo resultante}
\phantomsection\label{ejemplo-representaciuxf3n-del-flujo-resultante}
\begin{itemize}
\tightlist
\item
  El flujo resultante
  \(\mathbf{f} = (2.4;\, 5;\, 8;\, 8;\, 2.4;\, 2.6;\, 3)^T\) satisface
  la conservación en todos los nodos y constituye un flujo factible bien
  definido.
\end{itemize}
\end{frame}

\begin{frame}{Ejemplo: Combinación lineal de vectores de flujo}
\phantomsection\label{ejemplo-combinaciuxf3n-lineal-de-vectores-de-flujo}
Considérense los tres vectores de flujo
\(\mathbf{f}^1, \mathbf{f}^2, \mathbf{f}^3 \in \mathbb{R}^7\). Evaluamos
la combinación lineal:

\[\mathbf{f} = 2\mathbf{f}^1 + \mathbf{f}^2 - 3\mathbf{f}^3\]

\[
\begin{aligned}
\mathbf{f} &= 2(2.7;\, 2;\, 2;\, 2;\, 2.7;\, -0.7;\, 0)^T + (0;\, 1;\, 1;\, 1;\, 0;\, 1;\, 0)^T - 3(1;\, 0;\, -1;\, -1;\, 1;\, -1;\, -1)^T \\
&= (2.4;\, 5;\, 8;\, 8;\, 2.4;\, 2.6;\, 3)^T
\end{aligned}
\]
\end{frame}

\begin{frame}{Ejemplo: Verificación de conservación en los vértices}
\phantomsection\label{ejemplo-verificaciuxf3n-de-conservaciuxf3n-en-los-vuxe9rtices}
\begin{itemize}
\tightlist
\item
  \textbf{Verificación nodo a nodo del flujo
  \(\mathbf{f} = (2.4;\, 5;\, 8;\, 8;\, 2.4;\, 2.6;\, 3)^T\)}:

  \begin{itemize}
  \tightlist
  \item
    \textbf{Nodo \(a\)}: Salen \(f_1 = 2.4\), entran
    \(f_5 = 2.4 \implies 2.4 = 2.4\).
  \item
    \textbf{Nodo \(b\)}: Salen \(f_2 = 5\), entran
    \(f_1 + f_6 = 2.4 + 2.6 = 5 \implies 5 = 5\).
  \item
    \textbf{Nodo \(c\)}: Salen \(f_3 = 8\), entran
    \(f_2 + f_7 = 5 + 3 = 8 \implies 8 = 8\).
  \item
    \textbf{Nodo \(d\)}: Salen \(f_4 = 8\), entran
    \(f_3 = 8 \implies 8 = 8\).
  \item
    \textbf{Nodo \(e\)}: Salen \(f_5 + f_6 + f_7 = 2.4 + 2.6 + 3 = 8\),
    entra \(f_4 = 8 \implies 8 = 8\).
  \end{itemize}
\end{itemize}
\end{frame}

\begin{frame}{Definición de ciclo y clasificación de arcos}
\phantomsection\label{definiciuxf3n-de-ciclo-y-clasificaciuxf3n-de-arcos}
\begin{itemize}
\tightlist
\item
  \textbf{Concepto de Ciclo}:

  \begin{itemize}
  \tightlist
  \item
    Cadena simple \(\mu\) (sin repetir arcos) cuyo nodo inicial y final
    coinciden.
  \item
    Al fijar un sentido de giro de referencia, los arcos del digrafo se
    clasifican en:
  \end{itemize}
\item
  \textbf{Clasificación de arcos}:

  \begin{itemize}
  \tightlist
  \item
    \textbf{Arcos positivos (\(u \in \mu^+\))}: Su orientación en el
    digrafo coincide con el sentido del ciclo.
  \item
    \textbf{Arcos negativos (\(u \in \mu^-\))}: Su orientación en el
    digrafo es opuesta al sentido del ciclo.
  \item
    \textbf{Arcos neutros (\(u \notin \mu\))}: Arcos fuera de la cadena
    del ciclo.
  \end{itemize}
\item
  \textbf{Inversión de giro}:

  \begin{itemize}
  \tightlist
  \item
    Invertir la orientación de referencia intercambia los arcos
    positivos y negativos (\(\mu^+ \leftrightarrow \mu^-\)).
  \end{itemize}
\end{itemize}
\end{frame}

\begin{frame}{Clasificación visual de arcos en un ciclo}
\phantomsection\label{clasificaciuxf3n-visual-de-arcos-en-un-ciclo}
\begin{itemize}
\tightlist
\item
  Al recorrer la red según el sentido de giro elegido:

  \begin{itemize}
  \tightlist
  \item
    Arcos directos o positivos (\(\mu^+\)): \(1, 4, 3\).
  \item
    Arcos opuestos o negativos (\(\mu^-\)): \(6, 7\).
  \end{itemize}
\end{itemize}
\end{frame}

\begin{frame}{Definición: Vector característico de un ciclo}
\phantomsection\label{definiciuxf3n-vector-caracteruxedstico-de-un-ciclo}
\begin{tcolorbox}[enhanced jigsaw, coltitle=black, toprule=.15mm, opacitybacktitle=0.6, toptitle=1mm, leftrule=.75mm, breakable, left=2mm, opacityback=0, colback=white, bottomtitle=1mm, titlerule=0mm, title=\textcolor{quarto-callout-note-color}{\faInfo}\hspace{0.5em}{Vector característico de un ciclo}, colbacktitle=quarto-callout-note-color!10!white, rightrule=.15mm, bottomrule=.15mm, arc=.35mm, colframe=quarto-callout-note-color-frame]

Un ciclo \(\mu\) se representa algebraicamente mediante un vector
característico
\(\mathbf{\mu} = (\mu_1, \mu_2, \dots, \mu_m)^T \in \mathbb{R}^m\),
donde la componente asociada a cada arco \(u \in U\) se define como:

\[
\mu_u = \begin{cases}
+1 & \text{si } u \in \mu^+ \\
-1 & \text{si } u \in \mu^- \\
0 & \text{si } u \notin \mu
\end{cases}
\]

\end{tcolorbox}

\begin{itemize}
\tightlist
\item
  Traduce la estructura combinatoria de un ciclo a un vector de entradas
  \(\{-1, 0, 1\}\) en \(\mathbb{R}^m\).
\end{itemize}
\end{frame}

\begin{frame}{Ejemplo: Digrafo de referencia para vectores de ciclos}
\phantomsection\label{ejemplo-digrafo-de-referencia-para-vectores-de-ciclos}
\begin{itemize}
\tightlist
\item
  Digrafo de 5 nodos y 7 arcos utilizado para la evaluación de los
  vectores característicos de ciclos
  \(\mathbf{\mu}_1 \dots \mathbf{\mu}_5\).
\end{itemize}
\end{frame}

\begin{frame}{Ejemplo: Vectores característicos de ciclos}
\phantomsection\label{ejemplo-vectores-caracteruxedsticos-de-ciclos}
Considérese el digrafo de referencia de 5 nodos \(s, b, c, d, t\) y 7
arcos. Se identifican los siguientes ciclos
\(\mathbf{\mu}_k \in \mathbb{R}^7\):

\begin{itemize}
\item
  \(\mathbf{\mu}_1 = (1, 1, 1, 0, 0, 0, 0)^T\): Ciclo orientado
  \(s \to b \to c \to s\) (arcos \(1, 2, 3\)).
\item
  \(\mathbf{\mu}_2 = (0, -1, 0, 1, 1, 0, 0)^T\): Ciclo
  \(b \to d \to c \to b\) (arcos \(4, 5\) y opuesto de \(2\)).
\item
  \(\mathbf{\mu}_3 = (1, 0, 1, 1, 1, 0, 0)^T\): Ciclo exterior
  \(s \to b \to d \to c \to s\).
\item
  \(\mathbf{\mu}_5 = (-1, -1, -1, 0, 0, 0, 0)^T\): Ciclo inverso al
  ciclo \(\mathbf{\mu}_1\).
\item
  \textbf{Relaciones lineales}:
  \[\mathbf{\mu}_1 = -\mathbf{\mu}_5, \qquad \mathbf{\mu}_3 = \mathbf{\mu}_1 + \mathbf{\mu}_2\]
\end{itemize}
\end{frame}

\begin{frame}{Representación gráfica: Ciclos como flujos}
\phantomsection\label{representaciuxf3n-gruxe1fica-ciclos-como-flujos}
\begin{itemize}
\tightlist
\item
  Los vectores \(\mathbf{f}^2 = (0;\, 1;\, 1;\, 1;\, 0;\, 1;\, 0)^T\) y
  \(\mathbf{f}^3 = (1;\, 0;\, -1;\, -1;\, 1;\, -1;\, -1)^T\) son flujos
  válidos generados por ciclos que satisfacen la conservación en todos
  los nodos.
\end{itemize}
\end{frame}

\begin{frame}{Proposición: Los ciclos son flujos}
\phantomsection\label{proposiciuxf3n-los-ciclos-son-flujos}
\begin{tcolorbox}[enhanced jigsaw, coltitle=black, toprule=.15mm, opacitybacktitle=0.6, toptitle=1mm, leftrule=.75mm, breakable, left=2mm, opacityback=0, colback=white, bottomtitle=1mm, titlerule=0mm, title=\textcolor{quarto-callout-note-color}{\faInfo}\hspace{0.5em}{Los ciclos son flujos}, colbacktitle=quarto-callout-note-color!10!white, rightrule=.15mm, bottomrule=.15mm, arc=.35mm, colframe=quarto-callout-note-color-frame]

El vector característico \(\mathbf{\mu}\) asociado a cualquier ciclo en
un digrafo es un flujo.

\end{tcolorbox}

\begin{itemize}
\tightlist
\item
  \textbf{Verificación de la conservación nodal}:

  \begin{itemize}
  \tightlist
  \item
    Para \(i \notin \mu\): Ningún arco incidente pertenece al ciclo
    \(\implies 0 = 0\).
  \item
    Para \(i \in \mu\): El ciclo entra por un arco \(u_{in}\) y sale por
    \(u_{out}\).
  \item
    Analizando la orientación respecto al giro (positivo o negativo), la
    contribución al balance entrante \(\sum f_{in}\) es igual a la del
    saliente \(\sum f_{out}\).
  \end{itemize}
\end{itemize}
\end{frame}

\begin{frame}{Corolario: Variación de flujo mediante ciclos}
\phantomsection\label{corolario-variaciuxf3n-de-flujo-mediante-ciclos}
\begin{tcolorbox}[enhanced jigsaw, coltitle=black, toprule=.15mm, opacitybacktitle=0.6, toptitle=1mm, leftrule=.75mm, breakable, left=2mm, opacityback=0, colback=white, bottomtitle=1mm, titlerule=0mm, title=\textcolor{quarto-callout-note-color}{\faInfo}\hspace{0.5em}{Variación de flujo mediante ciclos}, colbacktitle=quarto-callout-note-color!10!white, rightrule=.15mm, bottomrule=.15mm, arc=.35mm, colframe=quarto-callout-note-color-frame]

Dado un flujo \(\mathbf{f}\) en un digrafo \(G\) y un ciclo
\(\mathbf{\mu}\), el vector resultante de sumar a \(\mathbf{f}\) el
vector asociado al ciclo multiplicado por una constante
\(\alpha \in \mathbb{R}\)
(\(\mathbf{f}' = \mathbf{f} + \alpha \mathbf{\mu}\)) es también un flujo
en \(G\).

\end{tcolorbox}

\begin{itemize}
\tightlist
\item
  Permite modificar localmente el caudal de una red sin alterar el
  equilibrio de los nodos.
\end{itemize}
\end{frame}

\begin{frame}{Definición: Independencia de flujos}
\phantomsection\label{definiciuxf3n-independencia-de-flujos}
\begin{tcolorbox}[enhanced jigsaw, coltitle=black, toprule=.15mm, opacitybacktitle=0.6, toptitle=1mm, leftrule=.75mm, breakable, left=2mm, opacityback=0, colback=white, bottomtitle=1mm, titlerule=0mm, title=\textcolor{quarto-callout-note-color}{\faInfo}\hspace{0.5em}{Independencia de flujos}, colbacktitle=quarto-callout-note-color!10!white, rightrule=.15mm, bottomrule=.15mm, arc=.35mm, colframe=quarto-callout-note-color-frame]

Dado un digrafo \(G = (V, U)\), un conjunto de flujos
\(\{\mathbf{f}^1, \mathbf{f}^2, \dots, \mathbf{f}^k\}\) es linealmente
independiente si sus correspondientes vectores en \(\mathbb{R}^m\) son
linealmente independientes.

\end{tcolorbox}

\begin{itemize}
\tightlist
\item
  La independencia lineal garantiza que ninguno de los flujos de la base
  pueda construirse combinando los restantes.
\end{itemize}
\end{frame}

\begin{frame}{Teorema: Base fundamental de ciclos}
\phantomsection\label{teorema-base-fundamental-de-ciclos}
\begin{tcolorbox}[enhanced jigsaw, coltitle=black, toprule=.15mm, opacitybacktitle=0.6, toptitle=1mm, leftrule=.75mm, breakable, left=2mm, opacityback=0, colback=white, bottomtitle=1mm, titlerule=0mm, title=\textcolor{quarto-callout-important-color}{\faExclamation}\hspace{0.5em}{Base fundamental de ciclos}, colbacktitle=quarto-callout-important-color!10!white, rightrule=.15mm, bottomrule=.15mm, arc=.35mm, colframe=quarto-callout-important-color-frame]

Sea \(G = (V, U)\) un digrafo conexo con \(n\) vértices y \(m\) arcos, y
sea \(H = (V, T)\) un árbol soporte de \(G\) (formado por \(n - 1\)
arcos).

Al añadir a \(T\) cualquier arco fuera del árbol
\(u \in U \setminus T\), se forma un único ciclo elemental \(\mu^u\) en
\(G\) tal que \(\mu^u_u = 1\). El conjunto de ciclos
\(H = \{\mu^u \mid u \in U \setminus T\}\) constituye una \textbf{base
del espacio vectorial de flujos} del digrafo \(G\).

\end{tcolorbox}

\begin{itemize}
\tightlist
\item
  \textbf{Aclaración teórica}: El número de vectores de ciclo necesarios
  para generar cualquier circulación en el digrafo queda determinado
  exclusivamente por la cantidad de arcos fuera del árbol soporte
  (\(\nu(G) = m - n + 1\)).
\end{itemize}
\end{frame}

\begin{frame}{Definición: Número ciclomático}
\phantomsection\label{definiciuxf3n-nuxfamero-ciclomuxe1tico}
\begin{tcolorbox}[enhanced jigsaw, coltitle=black, toprule=.15mm, opacitybacktitle=0.6, toptitle=1mm, leftrule=.75mm, breakable, left=2mm, opacityback=0, colback=white, bottomtitle=1mm, titlerule=0mm, title=\textcolor{quarto-callout-note-color}{\faInfo}\hspace{0.5em}{Número ciclomático}, colbacktitle=quarto-callout-note-color!10!white, rightrule=.15mm, bottomrule=.15mm, arc=.35mm, colframe=quarto-callout-note-color-frame]

Se define el \textbf{número ciclomático} \(\nu(G)\) de un digrafo conexo
\(G = (V, U)\) como la dimensión del espacio vectorial de sus flujos:

\[
\nu(G) = m - n + 1
\]

\end{tcolorbox}

\begin{itemize}
\tightlist
\item
  Basta con identificar \(m - n + 1\) cuerdas fuera del árbol soporte
  para caracterizar unívocamente todo el espacio de circulaciones.
\end{itemize}
\end{frame}

\begin{frame}{Ejemplo: Base de ciclos - Planteamiento y Pasos}
\phantomsection\label{ejemplo-base-de-ciclos---planteamiento-y-pasos}
Considérese el digrafo de \(n = 5\) nodos y \(m = 7\) arcos. Su número
ciclomático es \(\nu(G) = 7 - 5 + 1 = 3\).

\begin{itemize}
\tightlist
\item
  \textbf{Paso 1: Selección del Árbol Soporte y Cuerdas}

  \begin{itemize}
  \tightlist
  \item
    Árbol soporte \(T = \{u_1, u_4, u_5, u_7\}\) (4 arcos en
    \textbf{azul}).
  \item
    Cuerdas o arcos no básicos \(U \setminus T = \{u_2, u_3, u_6\}\) (3
    arcos en \textbf{gris}).
  \end{itemize}
\item
  \textbf{Paso 2: Generación de los 3 Ciclos Fundamentales}

  \begin{enumerate}
  \tightlist
  \item
    Añadir
    \(u_2 = (b, c) \implies \mathbf{\mu}^2 = (1, 1, 0, 0, 1, 0, -1)^T\).
  \item
    Añadir
    \(u_3 = (c, d) \implies \mathbf{\mu}^3 = (0, 0, 1, 1, 0, 0, 1)^T\).
  \item
    Añadir
    \(u_6 = (e, b) \implies \mathbf{\mu}^6 = (-1, 0, 0, 0, -1, 1, 0)^T\).
  \end{enumerate}
\item
  \textbf{Paso 3: Descomposición Unívoca de un Flujo} Dado el flujo
  \(\mathbf{f} = (3, 4, 6, 6, 3, 1, 2)^T\), sus coordenadas en la base
  son \(f_{u_2} = 4, f_{u_3} = 6, f_{u_6} = 1\):
  \[\mathbf{f} = 4 \mathbf{\mu}^2 + 6 \mathbf{\mu}^3 + 1 \mathbf{\mu}^6\]
\end{itemize}
\end{frame}

\begin{frame}{Ejemplo: Animación interactiva con controles}
\phantomsection\label{ejemplo-animaciuxf3n-interactiva-con-controles}
\end{frame}

\begin{frame}{Definición: Cociclo (Corte)}
\phantomsection\label{definiciuxf3n-cociclo-corte}
\begin{tcolorbox}[enhanced jigsaw, coltitle=black, toprule=.15mm, opacitybacktitle=0.6, toptitle=1mm, leftrule=.75mm, breakable, left=2mm, opacityback=0, colback=white, bottomtitle=1mm, titlerule=0mm, title=\textcolor{quarto-callout-note-color}{\faInfo}\hspace{0.5em}{Cociclo (Corte)}, colbacktitle=quarto-callout-note-color!10!white, rightrule=.15mm, bottomrule=.15mm, arc=.35mm, colframe=quarto-callout-note-color-frame]

Dado un digrafo conexo \(G = (V, U)\) y un subconjunto propio no vacío
de vértices \(A \subset V\), se define el \textbf{cociclo} \(w(A)\) como
el subconjunto de arcos que tienen un extremo en \(A\) y el otro en
\(V \setminus A\):

\[
w(A) = w^+(A) \cup w^-(A)
\]

donde \(w^+(A)\) son los arcos salientes de \(A\) y \(w^-(A)\) los arcos
entrantes a \(A\).

\end{tcolorbox}
\end{frame}

\begin{frame}{Definición: Vector característico de un cociclo}
\phantomsection\label{definiciuxf3n-vector-caracteruxedstico-de-un-cociclo}
\begin{tcolorbox}[enhanced jigsaw, coltitle=black, toprule=.15mm, opacitybacktitle=0.6, toptitle=1mm, leftrule=.75mm, breakable, left=2mm, opacityback=0, colback=white, bottomtitle=1mm, titlerule=0mm, title=\textcolor{quarto-callout-note-color}{\faInfo}\hspace{0.5em}{Vector característico de un cociclo}, colbacktitle=quarto-callout-note-color!10!white, rightrule=.15mm, bottomrule=.15mm, arc=.35mm, colframe=quarto-callout-note-color-frame]

Un cociclo \(w(A)\) se representa algebraicamente mediante un vector
característico
\(\mathbf{\theta} = (\theta_1, \theta_2, \dots, \theta_m)^T \in \mathbb{R}^m\),
cuyas componentes se definen como:

\[
\theta_u = \begin{cases}
+1 & \text{si } u \in w^+(A) \\
-1 & \text{si } u \in w^-(A) \\
0 & \text{si } u \notin w(A)
\end{cases}
\]

\end{tcolorbox}
\end{frame}

\begin{frame}{Representación gráfica de un cociclo}
\phantomsection\label{representaciuxf3n-gruxe1fica-de-un-cociclo}
\begin{itemize}
\tightlist
\item
  Para el conjunto de nodos \(A = \{s, b, c\}\) (en \textbf{azul}), los
  arcos del cociclo \(w(A)\) (en \textbf{rojo teja}) son
  \(u_4, u_7 \in w^+(A)\) (\(\theta = +1\)) y \(u_5 \in w^-(A)\)
  (\(\theta = -1\)), resultando
  \(\mathbf{\theta} = (0, 0, 0, 1, -1, 0, 1)^T\).
\end{itemize}
\end{frame}

\begin{frame}{Ejemplo: Representación gráfica de los dos cociclos}
\phantomsection\label{ejemplo-representaciuxf3n-gruxe1fica-de-los-dos-cociclos}
\begin{itemize}
\tightlist
\item
  Comparativa visual de las particiones \(A_1 = \{a, b, e\}\)
  (izquierda) y \(A_2 = \{b, d\}\) (derecha) con sus respectivos
  vectores característicos \(\mathbf{\theta}^1\) y
  \(\mathbf{\theta}^2\).
\end{itemize}
\end{frame}

\begin{frame}{Ejemplo: Evaluación de vectores de cociclo (Parte 1)}
\phantomsection\label{ejemplo-evaluaciuxf3n-de-vectores-de-cociclo-parte-1}
Evaluación de dos particiones sobre la red de 5 nodos y 7 arcos:

\begin{enumerate}
\tightlist
\item
  \textbf{Subconjunto \(A_1 = \{a, b, e\}\)}:

  \begin{itemize}
  \tightlist
  \item
    Nodos \(a, b, e\) en la región izquierda (\(A_1\)) y \(c, d\) en la
    derecha (\(V \setminus A_1\)).
  \item
    Arcos internos a \(A_1\): \(u_1, u_5, u_6 \implies \theta = 0\).
  \item
    Arcos salientes \(w^+(A_1)\): \(u_2=(b,c)\) y
    \(u_7=(e,c) \implies \theta_{u_2} = 1, \theta_{u_7} = 1\).
  \item
    Arcos entrantes \(w^-(A_1)\):
    \(u_4=(d,e) \implies \theta_{u_4} = -1\).
  \item
    Arcos externos a \(A_1\): \(u_3=(c,d) \implies \theta_{u_3} = 0\).
  \item
    Vector característico de \(w(A_1)\):
    \[\mathbf{\theta}^1 = (0;\, 1;\, 0;\, -1;\, 0;\, 0;\, 1)^T\]
  \end{itemize}
\end{enumerate}
\end{frame}

\begin{frame}{Ejemplo: Evaluación de vectores de cociclo (Parte 2)}
\phantomsection\label{ejemplo-evaluaciuxf3n-de-vectores-de-cociclo-parte-2}
\begin{enumerate}
\setcounter{enumi}{1}
\tightlist
\item
  \textbf{Subconjunto \(A_2 = \{b, d\}\)}:

  \begin{itemize}
  \tightlist
  \item
    Nodos \(b, d\) en la región superior (\(A_2\)) y \(a, c, e\) en la
    inferior (\(V \setminus A_2\)).
  \item
    Arcos salientes \(w^+(A_2)\): \(u_2=(b,c)\) y
    \(u_4=(d,e) \implies \theta_{u_2} = 1, \theta_{u_4} = 1\).
  \item
    Arcos entrantes \(w^-(A_2)\): \(u_1=(a,b)\), \(u_3=(c,d)\) y
    \(u_6=(e,b) \implies \theta_{u_1} = -1, \theta_{u_3} = -1, \theta_{u_6} = -1\).
  \item
    Arcos no incidentes: \(u_5, u_7 \implies \theta = 0\).
  \item
    Vector característico de \(w(A_2)\):
    \[\mathbf{\theta}^2 = (-1;\, 1;\, -1;\, 1;\, 0;\, -1;\, 0)^T\]
  \end{itemize}
\end{enumerate}
\end{frame}

\begin{frame}{Ortogonalidad fundamental entre ciclos y cociclos}
\phantomsection\label{ortogonalidad-fundamental-entre-ciclos-y-cociclos}
\begin{itemize}
\item
  \textbf{Producto Escalar Nulo}: Para cualquier ciclo \(\mathbf{\mu}\)
  y cualquier cociclo \(\mathbf{\theta}\) sobre el mismo digrafo
  \(G = (V, U)\):

  \[
  \mathbf{\mu}^T \mathbf{\theta} = \sum_{u \in U} \mu_u \theta_u = 0
  \]
\item
  \textbf{Significado Físico}:

  \begin{itemize}
  \tightlist
  \item
    El caudal neto transportado por cualquier circulación cerrada a
    través de un corte topológico es idénticamente nulo.
  \end{itemize}
\end{itemize}
\end{frame}

\begin{frame}{Subespacios ortogonales complementarios}
\phantomsection\label{subespacios-ortogonales-complementarios}
\begin{itemize}
\item
  \textbf{Subespacios Ortogonales en \(\mathbb{R}^m\)}:

  \begin{itemize}
  \tightlist
  \item
    Espacio de Flujos \(\mathcal{C}\) (\(\dim = m - n + 1\)).
  \item
    Espacio de Cociclos \(\mathcal{C}^*\) (\(\dim = n - 1\)).
  \end{itemize}
\item
  \textbf{Suma directa ortogonal}:
  \[\mathbb{R}^m = \mathcal{C} \oplus \mathcal{C}^*\]
\item
  \textbf{Complementariedad dimensional}: \[(m - n + 1) + (n - 1) = m\]
\end{itemize}
\end{frame}

\begin{frame}{Lema: Coloración de arcos de Minty}
\phantomsection\label{lema-coloraciuxf3n-de-arcos-de-minty}
\begin{tcolorbox}[enhanced jigsaw, coltitle=black, toprule=.15mm, opacitybacktitle=0.6, toptitle=1mm, leftrule=.75mm, breakable, left=2mm, opacityback=0, colback=white, bottomtitle=1mm, titlerule=0mm, title=\textcolor{quarto-callout-note-color}{\faInfo}\hspace{0.5em}{Coloración de arcos de Minty}, colbacktitle=quarto-callout-note-color!10!white, rightrule=.15mm, bottomrule=.15mm, arc=.35mm, colframe=quarto-callout-note-color-frame]

Sea \(G = (V, U)\) un digrafo conexo cuyos arcos se encuentran
coloreados con cuatro opciones: \textbf{negro} (\(N\)), \textbf{verde}
(\(V\)), \textbf{rojo} (\(R\)) o \textbf{sin color} (\(0\)), existiendo
al menos un arco negro \(u_0 \in U\).

Entonces se cumple \textbf{una y solo una} de las dos aseveraciones:

\begin{enumerate}
\tightlist
\item
  El arco negro \(u_0\) pertenece a un \textbf{ciclo elemental}
  coloreado (arcos negros positivos, verdes negativos y rojos
  arbitrarios).
\item
  El arco negro \(u_0\) pertenece a un \textbf{cociclo elemental}
  coloreado (arcos negros en \(w^+(A)\), verdes en \(w^-(A)\) y sin
  color arbitrarios).
\end{enumerate}

\end{tcolorbox}
\end{frame}

\begin{frame}{Demostración constructiva de Minty (Algoritmo)}
\phantomsection\label{demostraciuxf3n-constructiva-de-minty-algoritmo}
\begin{itemize}
\tightlist
\item
  \textbf{Sea el arco negro \(u_0 = (t, s)\)}:

  \begin{itemize}
  \tightlist
  \item
    \textbf{Paso 1: Etiquetado inicial}: Asignar \(p(s) = s\). El resto
    \(p(j) = 0\).
  \item
    \textbf{Paso 2: Propagación de etiquetas}: Dado \(i \in A\)
    (\(p(i) \neq 0\)) y \(j \notin A\) (\(p(j) = 0\)), asignar
    \(p(j) = i\) si:

    \begin{itemize}
    \tightlist
    \item
      Existe el arco saliente \((i, j)\) coloreado de \textbf{negro} o
      \textbf{rojo}.
    \item
      Existe el arco entrante \((j, i)\) coloreado de \textbf{verde} o
      \textbf{rojo}.
    \end{itemize}
  \item
    \textbf{Paso 3: Criterio de parada}:

    \begin{itemize}
    \tightlist
    \item
      \textbf{Caso A (\(p(t) \neq 0\))}: Se alcanza \(t \implies\)
      \textbf{Ciclo elemental} de \(t\) a \(s\) con \(u_0\).
    \item
      \textbf{Caso B (\(p(t) = 0\))}: No se alcanza \(t \implies\)
      \textbf{Cociclo elemental} sobre \(A = \{i \mid p(i) \neq 0\}\).
    \end{itemize}
  \end{itemize}
\end{itemize}
\end{frame}

\begin{frame}{Ejemplo Minty: Representación gráfica}
\phantomsection\label{ejemplo-minty-representaciuxf3n-gruxe1fica}
\begin{itemize}
\tightlist
\item
  Izquierda: Propagación de marcas hasta alcanzar el nodo origen \(b\)
  (Caso A - Ciclo).
\item
  Derecha: Propagación detenida sin alcanzar \(b\) (Caso B - Cociclo).
\end{itemize}
\end{frame}

\begin{frame}{Ejemplo Minty - Caso A y Caso B}
\phantomsection\label{ejemplo-minty---caso-a-y-caso-b}
\begin{itemize}
\tightlist
\item
  \textbf{Caso A: Generación de un Ciclo} (Gráfico Izquierdo)

  \begin{itemize}
  \tightlist
  \item
    Arco negro \(u_0 = (b, c)\). Arco \((e, b)\) en \textbf{rojo};
    \((e, c), (d, c), (d, e), (a, e)\) en \textbf{verde}.
  \item
    Propagación: \(p(c)=c \to p(e)=c \to p(d)=c \to p(b)=e\).
  \item
    Se alcanza \(b \implies\) Ciclo
    \(\mu = \{(e, b), (e, c), (b, c)\}\).
  \end{itemize}
\item
  \textbf{Caso B: Generación de un Cociclo} (Gráfico Derecho)

  \begin{itemize}
  \tightlist
  \item
    Arco negro \(u_0 = (b, c)\). Arcos \((e, b), (e, c), (c, d)\) en
    \textbf{verde}; \((d, e), (a, b)\) en \textbf{negro}.
  \item
    Propagación: \(p(c)=c \to p(e)=c\). Los arcos salientes verdes no
    admiten marca \(\implies p(b)=0\).
  \item
    Cociclo \(w(A)\) sobre \(A = \{c, e\}\) con \(u_0 \in w^-(A)\).
  \end{itemize}
\end{itemize}
\end{frame}

\begin{frame}{Observación: Elementalidad de los ciclos y cociclos}
\phantomsection\label{observaciuxf3n-elementalidad-de-los-ciclos-y-cociclos}
\begin{tcolorbox}[enhanced jigsaw, coltitle=black, toprule=.15mm, opacitybacktitle=0.6, toptitle=1mm, leftrule=.75mm, breakable, left=2mm, opacityback=0, colback=white, bottomtitle=1mm, titlerule=0mm, title=\textcolor{quarto-callout-note-color}{\faInfo}\hspace{0.5em}{Elementalidad de los ciclos y cociclos de Minty}, colbacktitle=quarto-callout-note-color!10!white, rightrule=.15mm, bottomrule=.15mm, arc=.35mm, colframe=quarto-callout-note-color-frame]

\begin{enumerate}
\item
  \textbf{Elementalidad del Ciclo}: Al alcanzar el origen \(t\), la
  asignación de antecesores \(p(j)\) asigna un único padre a cada nodo,
  impidiendo vértices duplicados en la trayectoria de retorno.
\item
  \textbf{Elementalidad del Cociclo}: La partición \(A\) (etiquetados) y
  \(V \setminus A\) (no etiquetados) genera subgrafos inducidos \(G[A]\)
  y \(G[V \setminus A]\) conexos por construcción inductiva del
  etiquetado.
\end{enumerate}

\end{tcolorbox}
\end{frame}

\section{El problema del flujo
máximo}\label{el-problema-del-flujo-muxe1ximo}

\begin{frame}{Introducción al problema del flujo máximo}
\phantomsection\label{introducciuxf3n-al-problema-del-flujo-muxe1ximo}
\begin{itemize}
\tightlist
\item
  \textbf{Motivación estratégica}:

  \begin{itemize}
  \tightlist
  \item
    Determinar el volumen máximo de mercancía, datos o fluidos que puede
    enviarse desde un nodo fuente \(s\) a un nodo sumidero \(t\).
  \item
    Respetar la capacidad física máxima \(c_u\) de cada conexión.
  \end{itemize}
\item
  \textbf{Relevancia práctica}:

  \begin{itemize}
  \tightlist
  \item
    Logística de distribución y transporte.
  \item
    Capacidad de tráfico en redes de comunicaciones e infraestructuras.
  \end{itemize}
\end{itemize}
\end{frame}

\begin{frame}{Definición: Red de transporte}
\phantomsection\label{definiciuxf3n-red-de-transporte}
\begin{tcolorbox}[enhanced jigsaw, coltitle=black, toprule=.15mm, opacitybacktitle=0.6, toptitle=1mm, leftrule=.75mm, breakable, left=2mm, opacityback=0, colback=white, bottomtitle=1mm, titlerule=0mm, title=\textcolor{quarto-callout-note-color}{\faInfo}\hspace{0.5em}{Red de transporte}, colbacktitle=quarto-callout-note-color!10!white, rightrule=.15mm, bottomrule=.15mm, arc=.35mm, colframe=quarto-callout-note-color-frame]

Una \textbf{red de transporte} es un par \((G = (V, U), \mathbf{c})\),
donde \(G = (V, U)\) es un digrafo conexo sin bucles y
\(\mathbf{c} = \{c_u \mid u \in U\}\) representa el conjunto de
\textbf{capacidades} superiores no negativas asociadas a cada arco
\(u \in U\) (\(c_u \ge 0\)).

\end{tcolorbox}

\begin{itemize}
\tightlist
\item
  La capacidad \(c_u\) fija la cota física máxima de caudal por el arco
  \(u\).
\end{itemize}
\end{frame}

\begin{frame}{Definición: Flujo compatible (preliminar)}
\phantomsection\label{definiciuxf3n-flujo-compatible-preliminar}
\begin{tcolorbox}[enhanced jigsaw, coltitle=black, toprule=.15mm, opacitybacktitle=0.6, toptitle=1mm, leftrule=.75mm, breakable, left=2mm, opacityback=0, colback=white, bottomtitle=1mm, titlerule=0mm, title=\textcolor{quarto-callout-note-color}{\faInfo}\hspace{0.5em}{Flujo compatible (versión preliminar)}, colbacktitle=quarto-callout-note-color!10!white, rightrule=.15mm, bottomrule=.15mm, arc=.35mm, colframe=quarto-callout-note-color-frame]

Dada una red de transporte \((G, \mathbf{c})\), se define un
\textbf{flujo compatible} como un vector \(\mathbf{f} \in \mathbb{R}^m\)
que satisface:

\begin{enumerate}
\tightlist
\item
  \textbf{Ley de conservación pura}:
  \(\sum_{u \in w^+(i)} f_u = \sum_{u \in w^-(i)} f_u, \quad \forall i \in V\).
\item
  \textbf{Compatibilidad con capacidades}:
  \(0 \le f_u \le c_u, \quad \forall u \in U\).
\end{enumerate}

\end{tcolorbox}
\end{frame}

\begin{frame}{Paradoja de la conservación nula en s y t}
\phantomsection\label{paradoja-de-la-conservaciuxf3n-nula-en-s-y-t}
\begin{itemize}
\tightlist
\item
  \textbf{Conflicto en la fuente \(s\) y el sumidero \(t\)}:

  \begin{itemize}
  \tightlist
  \item
    De la fuente \(s\) solo salen arcos (\(w^-(s) = \emptyset\)). Si
    exigimos conservación pura,
    \(\sum_{entrantes} f_u = 0 \implies \sum_{salientes} f_u = 0 \implies f_u = 0, \forall u \in w^+(s)\).
  \item
    Al sumidero \(t\) solo entran arcos (\(w^+(t) = \emptyset\)). Exigir
    conservación pura impone \(\sum_{entrantes} f_u = 0\).
  \end{itemize}
\item
  \textbf{Efecto de bloqueo}:

  \begin{itemize}
  \tightlist
  \item
    Exigir conservación estricta en todos los nodos sin excepción fuerza
    a que el \textbf{único flujo compatible sea el flujo nulo}
    (\(\mathbf{f} = \mathbf{0}\)).
  \end{itemize}
\end{itemize}
\end{frame}

\begin{frame}{Ilustración: Paradoja del flujo nulo}
\phantomsection\label{ilustraciuxf3n-paradoja-del-flujo-nulo}
\begin{itemize}
\tightlist
\item
  Si se impone la ley de conservación pura en \(s\) y \(t\), la única
  solución admisible es \(\mathbf{f} = \mathbf{0}\), impidiendo el envío
  de flujo útil.
\end{itemize}
\end{frame}

\begin{frame}{Concepto: Arco ficticio de retorno}
\phantomsection\label{concepto-arco-ficticio-de-retorno}
\begin{itemize}
\tightlist
\item
  \textbf{Construcción del arco ficticio de retorno \(u_0 = (t, s)\)}:

  \begin{itemize}
  \tightlist
  \item
    Se conecta el sumidero \(t\) de regreso a la fuente \(s\) mediante
    \(u_0 = (t, s)\) con capacidad ilimitada (\(c_0 = \infty\)).
  \item
    Todo el flujo que llega a \(t\) recircula de vuelta a \(s\),
    convirtiendo el transporte en una \textbf{circulación cerrada}.
  \end{itemize}
\item
  \textbf{Valor del flujo \(v(\mathbf{f})\)}:

  \begin{itemize}
  \tightlist
  \item
    El caudal por el arco ficticio representa el flujo neto
    transportado:
    \[f_0 = v(\mathbf{f}) = \sum_{j \in \Gamma^+(s)} f_{sj} = \sum_{j \in \Gamma^-(t)} f_{jt}\]
  \end{itemize}
\end{itemize}
\end{frame}

\begin{frame}{Definición: (s,t)-Flujo compatible}
\phantomsection\label{definiciuxf3n-st-flujo-compatible}
\begin{tcolorbox}[enhanced jigsaw, coltitle=black, toprule=.15mm, opacitybacktitle=0.6, toptitle=1mm, leftrule=.75mm, breakable, left=2mm, opacityback=0, colback=white, bottomtitle=1mm, titlerule=0mm, title=\textcolor{quarto-callout-note-color}{\faInfo}\hspace{0.5em}{(s, t)-Flujo compatible}, colbacktitle=quarto-callout-note-color!10!white, rightrule=.15mm, bottomrule=.15mm, arc=.35mm, colframe=quarto-callout-note-color-frame]

Dada una red de transporte \((G = (V, U), \mathbf{c})\) y dos vértices
\(s, t \in V\), un \textbf{\((s, t)\)-flujo compatible} es una
circulación factible sobre la red ampliada
\(G' = (V, U \cup \{u_0=(t,s)\})\) con \(c_0 = \infty\), tal que
\(0 \le f_u \le c_u, \forall u \in U\).

\end{tcolorbox}

\begin{itemize}
\tightlist
\item
  Permite formular la conservación en \textbf{todos} los nodos
  incluyendo \(s\) y \(t\).
\end{itemize}
\end{frame}

\begin{frame}{Ilustración: Red ampliada con retorno ficticio}
\phantomsection\label{ilustraciuxf3n-red-ampliada-con-retorno-ficticio}
\begin{itemize}
\tightlist
\item
  El arco ficticio \(u_0 = (t, s)\) lleva un valor de flujo \(f_0 = 8\),
  permitiendo cerrar la red y definir \(v(\mathbf{f}) = 8\).
\end{itemize}
\end{frame}

\begin{frame}{Formulación del problema del flujo máximo}
\phantomsection\label{formulaciuxf3n-del-problema-del-flujo-muxe1ximo}
\begin{tcolorbox}[enhanced jigsaw, coltitle=black, toprule=.15mm, opacitybacktitle=0.6, toptitle=1mm, leftrule=.75mm, breakable, left=2mm, opacityback=0, colback=white, bottomtitle=1mm, titlerule=0mm, title=\textcolor{quarto-callout-note-color}{\faInfo}\hspace{0.5em}{Problema del flujo máximo}, colbacktitle=quarto-callout-note-color!10!white, rightrule=.15mm, bottomrule=.15mm, arc=.35mm, colframe=quarto-callout-note-color-frame]

Dada una red de transporte \((G = (V, U), \mathbf{c})\) y dos vértices
\(s, t \in V\), el \textbf{problema del flujo máximo} consiste en
determinar un \((s, t)\)-flujo compatible
\(\mathbf{f}^* \in \mathbb{R}^m\) que maximice el valor \(f_0\) del arco
de retorno \(u_0 = (t, s)\):

\[
\begin{aligned}
\max_{\mathbf{f} \in \mathbb{R}^m} \quad & f_0 = \sum_{j \in \Gamma^+(s)} f_{sj} = \sum_{j \in \Gamma^-(t)} f_{jt} \\
\text{s.a.} \quad & \sum_{u \in w^+(i)} f_u - \sum_{u \in w^-(i)} f_u = 0, \quad \forall i \in V \\
& 0 \le f_u \le c_u, \quad \forall u \in U
\end{aligned}
\]

\end{tcolorbox}
\end{frame}

\begin{frame}{Ejemplo: Evaluación gráfica de compatibilidad}
\phantomsection\label{ejemplo-evaluaciuxf3n-gruxe1fica-de-compatibilidad}
\begin{itemize}
\tightlist
\item
  Evaluación de la compatibilidad de caudales sobre una red de 6 nodos
  con un valor de flujo \(f_0 = 8\).
\end{itemize}
\end{frame}

\begin{frame}{Ejemplo: Compatibilidad y saturación (Análisis)}
\phantomsection\label{ejemplo-compatibilidad-y-saturaciuxf3n-anuxe1lisis}
Considérese la red de 6 nodos \(V = \{s, b, c, d, e, t\}\) con flujo de
retorno \(f_0 = 8\):

\begin{itemize}
\tightlist
\item
  \textbf{Flujo compatible base}: Todos los arcos satisfacen
  \(0 \le f_u \le c_u\).

  \begin{itemize}
  \tightlist
  \item
    Arco \((s, b)\): \(f_{sb} = 5 \le c_{sb} = 5 \implies\) \textbf{Arco
    saturado}.
  \item
    Arco \((b, c)\): \(f_{bc} = 2 \le c_{bc} = 3 \implies\) \textbf{Arco
    no saturado}.
  \end{itemize}
\item
  \textbf{Intento de modificación por ciclo}: Añadir 2 unidades por el
  ciclo \(c \to d \to e \to c\):

  \begin{itemize}
  \tightlist
  \item
    Arco \((c, d)\): \(f_{cd} = 1 + 2 = 3 > c_{cd} = 1 \implies\)
    \textbf{Violación de capacidad}.
  \item
    Arco \((d, e)\): \(f_{de} = 4 + 2 = 6 > c_{de} = 3 \implies\)
    \textbf{Violación de capacidad}.
  \end{itemize}
\item
  \textbf{Conclusión}: El incremento por este ciclo destruye la
  compatibilidad del flujo.
\end{itemize}
\end{frame}

\begin{frame}{Ejemplo: Red de tráfico urbano}
\phantomsection\label{ejemplo-red-de-truxe1fico-urbano}
\begin{itemize}
\tightlist
\item
  \textbf{Planteamiento}: Red vial con intersecciones
  \(V = \{s, a, b, t\}\) y capacidades \(c_u\) (vehículos por minuto):

  \begin{itemize}
  \tightlist
  \item
    \(c_{sa} = 100\), \(c_{sb} = 100\), \(c_{ab} = 50\),
    \(c_{at} = 75\), \(c_{bt} = 100\).
  \end{itemize}
\end{itemize}
\end{frame}

\begin{frame}{Ejemplo: Resolución del tráfico urbano}
\phantomsection\label{ejemplo-resoluciuxf3n-del-truxe1fico-urbano}
\begin{itemize}
\tightlist
\item
  \textbf{Análisis de cuellos de botella}:

  \begin{itemize}
  \tightlist
  \item
    Inyección desde la fuente \(s\):
    \(c_{sa} + c_{sb} = 100 + 100 = 200\) veh/min.
  \item
    Capacidad de llegada al sumidero \(t\):
    \(c_{at} + c_{bt} = 75 + 100 = 175\) veh/min.
  \end{itemize}
\item
  \textbf{Conclusión}:

  \begin{itemize}
  \tightlist
  \item
    Aunque \(s\) puede emitir 200 veh/min, las vías de acceso a \(t\)
    restringen el paso a un máximo de 175 veh/min.
  \item
    Valor del flujo máximo: \(v(\mathbf{f}^*) = 175\) veh/min.
  \end{itemize}
\end{itemize}
\end{frame}

\begin{frame}{Ejemplo: Asignación de tareas a empleados}
\phantomsection\label{ejemplo-asignaciuxf3n-de-tareas-a-empleados}
\begin{itemize}
\tightlist
\item
  \textbf{Planteamiento}: 3 empleados (\(E_1, E_2, E_3\)) y 3 tareas
  (\(T_1, T_2, T_3\)) con arcos de compatibilidad.
\end{itemize}
\end{frame}

\begin{frame}{Ejemplo: Asignación de tareas (Transformación)}
\phantomsection\label{ejemplo-asignaciuxf3n-de-tareas-transformaciuxf3n}
\begin{itemize}
\tightlist
\item
  \textbf{Transformación a Red de Transporte}:

  \begin{enumerate}
  \tightlist
  \item
    Nodo fuente \(s\) conectado a cada empleado \(E_i\) con
    \(c_{s E_i} = 1\) (1 tarea por empleado).
  \item
    Arcos de compatibilidad \((E_i, T_j)\) con capacidad
    \(c_{E_i T_j} = 1\).
  \item
    Nodo sumidero \(t\) conectado desde cada tarea \(T_j\) con
    \(c_{T_j t} = 1\) (1 empleado por tarea).
  \item
    Arco ficticio de retorno \(u_0 = (t, s)\) con \(c_0 = \infty\).
  \end{enumerate}
\item
  Resolver el flujo máximo sobre la red equivale a hallar el número
  máximo de tareas ejecutables.
\end{itemize}
\end{frame}

\begin{frame}{Esquema iterativo de aumento de flujo}
\phantomsection\label{esquema-iterativo-de-aumento-de-flujo}
\begin{itemize}
\tightlist
\item
  \textbf{Procedimiento iterativo general}:

  \begin{enumerate}
  \tightlist
  \item
    \textbf{Flujo Inicial}: Comenzar con un flujo compatible válido
    \(\mathbf{f}^0\) sobre la red (típicamente el flujo trivial nulo
    \(\mathbf{f}^0 = \mathbf{0}\)).
  \item
    \textbf{Búsqueda de Aumento}: Buscar un nuevo flujo compatible
    \(\mathbf{f}^{k+1}\) que incremente el valor circulante por el arco
    de retorno \(f_0\).
  \item
    \textbf{Iteración y Criterio de Parada}: Repetir el procedimiento a
    partir del nuevo flujo hasta que no sea posible realizar ningún
    incremento adicional.
  \end{enumerate}
\item
  \textbf{Justificación algebraica}:

  \begin{itemize}
  \tightlist
  \item
    Toda diferencia entre flujos \(\mathbf{f}^2 - \mathbf{f}^1\) es una
    \textbf{combinación lineal de ciclos}.
  \item
    Para aumentar \(f_0\), el algoritmo busca un ciclo que incluya al
    arco ficticio \(u_0 = (t, s)\) y admita un incremento de caudal.
  \end{itemize}
\end{itemize}
\end{frame}

\begin{frame}{Definición: Cadena incremental y ciclo incremental}
\phantomsection\label{definiciuxf3n-cadena-incremental-y-ciclo-incremental}
\begin{tcolorbox}[enhanced jigsaw, coltitle=black, toprule=.15mm, opacitybacktitle=0.6, toptitle=1mm, leftrule=.75mm, breakable, left=2mm, opacityback=0, colback=white, bottomtitle=1mm, titlerule=0mm, title=\textcolor{quarto-callout-note-color}{\faInfo}\hspace{0.5em}{Cadena incremental}, colbacktitle=quarto-callout-note-color!10!white, rightrule=.15mm, bottomrule=.15mm, arc=.35mm, colframe=quarto-callout-note-color-frame]

Dado un flujo compatible \(\mathbf{f}\) en la red \((G, \mathbf{c})\),
una cadena \(\mu[a, b]\) entre \(a\) y \(b\) es una \textbf{cadena
incremental} (o camino aumentante) si satisface:

\[
f_u < c_u, \quad \forall u \in \mu^+ \qquad \text{y} \qquad f_u > 0, \quad \forall u \in \mu^-
\]

\end{tcolorbox}

\begin{tcolorbox}[enhanced jigsaw, coltitle=black, toprule=.15mm, opacitybacktitle=0.6, toptitle=1mm, leftrule=.75mm, breakable, left=2mm, opacityback=0, colback=white, bottomtitle=1mm, titlerule=0mm, title=\textcolor{quarto-callout-note-color}{\faInfo}\hspace{0.5em}{Ciclo incremental}, colbacktitle=quarto-callout-note-color!10!white, rightrule=.15mm, bottomrule=.15mm, arc=.35mm, colframe=quarto-callout-note-color-frame]

Un \textbf{ciclo incremental} es una cadena incremental cerrada. En el
flujo máximo, se busca una cadena incremental entre \(s\) y \(t\) que
complete con \(u_0 = (t, s)\) un ciclo incremental.

\end{tcolorbox}
\end{frame}

\begin{frame}{Definición: Capacidad residual}
\phantomsection\label{definiciuxf3n-capacidad-residual}
\begin{tcolorbox}[enhanced jigsaw, coltitle=black, toprule=.15mm, opacitybacktitle=0.6, toptitle=1mm, leftrule=.75mm, breakable, left=2mm, opacityback=0, colback=white, bottomtitle=1mm, titlerule=0mm, title=\textcolor{quarto-callout-note-color}{\faInfo}\hspace{0.5em}{Capacidad residual de una cadena}, colbacktitle=quarto-callout-note-color!10!white, rightrule=.15mm, bottomrule=.15mm, arc=.35mm, colframe=quarto-callout-note-color-frame]

La \textbf{capacidad residual} \(c(\mu)\) de una cadena incremental
\(\mu[s, t]\) es el máximo flujo adicional que puede enviarse a través
de ella:

\[
c(\mu) = \min \left\{ \min_{u \in \mu^+} (c_u - f_u), \, \min_{u \in \mu^-} f_u \right\}
\]

\end{tcolorbox}

\begin{itemize}
\tightlist
\item
  Aumentando el flujo en \(+c(\mu)\) en arcos directos \(u \in \mu^+\) y
  reduciéndolo en \(-c(\mu)\) en arcos inversos \(u \in \mu^-\), el
  valor total \(f_0\) aumenta exactamente en \(c(\mu)\) unidades
  manteniendo la compatibilidad estricta.
\end{itemize}
\end{frame}

\begin{frame}{Flujo máximo en grafos no dirigidos: Planteamiento}
\phantomsection\label{flujo-muxe1ximo-en-grafos-no-dirigidos-planteamiento}
\begin{itemize}
\tightlist
\item
  \textbf{Tránsito bidireccional en aristas no dirigidas}:

  \begin{itemize}
  \tightlist
  \item
    En redes sobre un grafo no dirigido \(G = (V, E)\), cada arista
    \(e = \{a, b\}\) con capacidad \(c_e\) admite flujo en ambas
    direcciones: \((a, b)\) y \((b, a)\).
  \item
    \textbf{Restricción de capacidad global}:
    \[f_{(a,b)} + f_{(b,a)} \le c_e, \qquad f_{(a,b)} \ge 0, \quad f_{(b,a)} \ge 0\]
  \end{itemize}
\item
  \textbf{Necesidad de transformación}:

  \begin{itemize}
  \tightlist
  \item
    Para aplicar algoritmos de flujo sobre digrafos, se debe reemplazar
    cada arista no dirigida por un \textbf{gadget dirigido} que impida
    el tránsito simultáneo en ambos sentidos.
  \end{itemize}
\end{itemize}
\end{frame}

\begin{frame}{Definición: Gadget de 5 arcos dirigidos}
\phantomsection\label{definiciuxf3n-gadget-de-5-arcos-dirigidos}
\begin{tcolorbox}[enhanced jigsaw, coltitle=black, toprule=.15mm, opacitybacktitle=0.6, toptitle=1mm, leftrule=.75mm, breakable, left=2mm, opacityback=0, colback=white, bottomtitle=1mm, titlerule=0mm, title=\textcolor{quarto-callout-note-color}{\faInfo}\hspace{0.5em}{Transformación de aristas no dirigidas (Gadget de 5 arcos)}, colbacktitle=quarto-callout-note-color!10!white, rightrule=.15mm, bottomrule=.15mm, arc=.35mm, colframe=quarto-callout-note-color-frame]

Cada arista no dirigida \(e = \{a, b\}\) con capacidad \(c_e = c\) se
sustituye por un conjunto de \textbf{cinco arcos dirigidos} y dos
vértices auxiliares \(a', b'\):

\[
U_e = \{ (a, a'), \, (b, a'), \, (a', b'), \, (b', a), \, (b', b) \}
\]

asignando capacidad \(c\) a todos los arcos del conjunto \(U_e\).

\end{tcolorbox}
\end{frame}

\begin{frame}{Ilustración: Gadget de 5 arcos}
\phantomsection\label{ilustraciuxf3n-gadget-de-5-arcos}
\begin{itemize}
\tightlist
\item
  Transformación de una arista no dirigida \(\{a, b\}\) de capacidad
  \(c\) (izquierda) en un gadget equivalente de 5 arcos dirigidos con
  nodos auxiliares \(a', b'\) (derecha).
\end{itemize}
\end{frame}

\begin{frame}{Funcionamiento y equivalencia del gadget}
\phantomsection\label{funcionamiento-y-equivalencia-del-gadget}
\begin{itemize}
\tightlist
\item
  \textbf{Tránsito de \(a\) hacia \(b\)}:

  \begin{itemize}
  \tightlist
  \item
    El flujo sigue el camino \(a \to a' \to b' \to b\), atravesando el
    arco central \((a', b')\) con capacidad \(c\).
  \end{itemize}
\item
  \textbf{Tránsito de \(b\) hacia \(a\)}:

  \begin{itemize}
  \tightlist
  \item
    El flujo sigue el camino \(b \to a' \to b' \to a\), atravesando el
    \textbf{mismo} arco central \((a', b')\).
  \end{itemize}
\item
  \textbf{Garantía de compatibilidad}:

  \begin{itemize}
  \tightlist
  \item
    Como ambas direcciones deben atravesar obligatoriamente el arco
    central \((a', b')\), se satisface de forma automática
    \(f_{(a,b)} + f_{(b,a)} \le c\), impidiendo la circulación
    simultánea contrapuesta.
  \end{itemize}
\end{itemize}
\end{frame}

\begin{frame}{Flujo máximo y corte mínimo: Planteamiento}
\phantomsection\label{flujo-muxe1ximo-y-corte-muxednimo-planteamiento}
\begin{itemize}
\tightlist
\item
  \textbf{Conexión entre flujo y barreras topológicas}:

  \begin{itemize}
  \tightlist
  \item
    El flujo máximo entre \(s\) y \(t\) no está restringido solo por
    capacidades individuales, sino por barreras que separan el origen
    del destino.
  \end{itemize}
\item
  \textbf{Concepto de corte}:

  \begin{itemize}
  \tightlist
  \item
    Partición de los nodos en dos conjuntos que aísla la fuente \(s\) en
    un lado y el sumidero \(t\) en el otro.
  \end{itemize}
\end{itemize}
\end{frame}

\begin{frame}{Definición: (s,t)-Corte}
\phantomsection\label{definiciuxf3n-st-corte}
\begin{tcolorbox}[enhanced jigsaw, coltitle=black, toprule=.15mm, opacitybacktitle=0.6, toptitle=1mm, leftrule=.75mm, breakable, left=2mm, opacityback=0, colback=white, bottomtitle=1mm, titlerule=0mm, title=\textcolor{quarto-callout-note-color}{\faInfo}\hspace{0.5em}{(s, t)-Corte}, colbacktitle=quarto-callout-note-color!10!white, rightrule=.15mm, bottomrule=.15mm, arc=.35mm, colframe=quarto-callout-note-color-frame]

Dada la red de transporte \((G = (V, U), \mathbf{c})\), un \textbf{corte
que separa \(s\) y \(t\)} o \textbf{\((s, t)\)-corte} es el conjunto de
arcos salientes \(w^+(N)\) asociado a un subconjunto de vértices
\(N \subset V\) tal que \(s \in N\) y \(t \notin N\).

\end{tcolorbox}
\end{frame}

\begin{frame}{Definición: Capacidad de un corte}
\phantomsection\label{definiciuxf3n-capacidad-de-un-corte}
\begin{tcolorbox}[enhanced jigsaw, coltitle=black, toprule=.15mm, opacitybacktitle=0.6, toptitle=1mm, leftrule=.75mm, breakable, left=2mm, opacityback=0, colback=white, bottomtitle=1mm, titlerule=0mm, title=\textcolor{quarto-callout-note-color}{\faInfo}\hspace{0.5em}{Capacidad de un corte}, colbacktitle=quarto-callout-note-color!10!white, rightrule=.15mm, bottomrule=.15mm, arc=.35mm, colframe=quarto-callout-note-color-frame]

Dada la red de transporte \((G = (V, U), \mathbf{c})\), se define la
\textbf{capacidad de un \((s, t)\)-corte} como la suma de las
capacidades de los arcos salientes que lo componen:

\[
c(w^+(N)) = \sum_{u \in w^+(N)} c_u
\]

\end{tcolorbox}
\end{frame}

\begin{frame}{Definición: Problema de corte de mínima capacidad}
\phantomsection\label{definiciuxf3n-problema-de-corte-de-muxednima-capacidad}
\begin{tcolorbox}[enhanced jigsaw, coltitle=black, toprule=.15mm, opacitybacktitle=0.6, toptitle=1mm, leftrule=.75mm, breakable, left=2mm, opacityback=0, colback=white, bottomtitle=1mm, titlerule=0mm, title=\textcolor{quarto-callout-note-color}{\faInfo}\hspace{0.5em}{Problema de corte de mínima capacidad}, colbacktitle=quarto-callout-note-color!10!white, rightrule=.15mm, bottomrule=.15mm, arc=.35mm, colframe=quarto-callout-note-color-frame]

El \textbf{problema de corte de mínima capacidad} consiste en determinar
la partición \(N \subset V\) (\(s \in N, t \notin N\)) que minimice la
capacidad del corte:

\[
\min_{\substack{N \subset V \\ s \in N, \, t \notin N}} c(w^+(N)) = \sum_{u \in w^+(N)} c_u
\]

\end{tcolorbox}
\end{frame}

\begin{frame}{Proposición: Acotación del flujo por la capacidad del
corte}
\phantomsection\label{proposiciuxf3n-acotaciuxf3n-del-flujo-por-la-capacidad-del-corte}
\begin{tcolorbox}[enhanced jigsaw, coltitle=black, toprule=.15mm, opacitybacktitle=0.6, toptitle=1mm, leftrule=.75mm, breakable, left=2mm, opacityback=0, colback=white, bottomtitle=1mm, titlerule=0mm, title=\textcolor{quarto-callout-note-color}{\faInfo}\hspace{0.5em}{Acotación del flujo por la capacidad del corte}, colbacktitle=quarto-callout-note-color!10!white, rightrule=.15mm, bottomrule=.15mm, arc=.35mm, colframe=quarto-callout-note-color-frame]

Para cualquier \((s, t)\)-flujo compatible \(\mathbf{f}\) y cualquier
\((s, t)\)-corte \(w^+(N)\), el valor del flujo está acotado
superiormente por la capacidad del corte:

\[
f_0 \le c(w^+(N)) = \sum_{u \in w^+(N)} c_u, \qquad \forall N \subset V \, (s \in N, t \notin N)
\]

\end{tcolorbox}
\end{frame}

\begin{frame}{Interpretación: Acotación por cuellos de botella}
\phantomsection\label{interpretaciuxf3n-acotaciuxf3n-por-cuellos-de-botella}
\begin{itemize}
\tightlist
\item
  \textbf{Interpretación física de la cota}:

  \begin{itemize}
  \tightlist
  \item
    Todo el flujo que se desplaza de \(s\) a \(t\) debe cruzar de forma
    obligatoria cualquier \((s, t)\)-corte.
  \item
    Ningún \((s, t)\)-flujo puede superar la capacidad del corte más
    estrecho:
    \[\max f_0 \le \min_{\substack{N \subset V \\ s \in N, \, t \notin N}} c(w^+(N))\]
  \end{itemize}
\end{itemize}
\end{frame}

\begin{frame}{Ejemplo: Evaluación gráfica de cortes}
\phantomsection\label{ejemplo-evaluaciuxf3n-gruxe1fica-de-cortes}
\begin{itemize}
\tightlist
\item
  Evaluación de tres \((s, t)\)-cortes sobre la red de 6 nodos con un
  flujo circulante de valor \(f_0 = 8\).
\end{itemize}
\end{frame}

\begin{frame}{Ejemplo: Evaluación de 3 cortes y capacidades}
\phantomsection\label{ejemplo-evaluaciuxf3n-de-3-cortes-y-capacidades}
Evaluación de los tres cortes sobre la red de 6 nodos:

\begin{enumerate}
\tightlist
\item
  \textbf{Corte 1 (\(N_1 = \{s\}\))} (línea roja teja):

  \begin{itemize}
  \tightlist
  \item
    \(w^+(N_1) = \{(s, b), (s, c)\} \implies c(w^+(N_1)) = 5 + 3 = 8\).
  \end{itemize}
\item
  \textbf{Corte 2 (\(N_2 = \{s, b, c\}\))} (línea azul cerúleo):

  \begin{itemize}
  \tightlist
  \item
    \(w^+(N_2) = \{(b, d), (c, d), (c, e)\} \implies c(w^+(N_2)) = 4 + 4 + 3 = 11\).
  \end{itemize}
\item
  \textbf{Corte 3 (\(N_3 = \{s, b, c, d, e\}\))} (línea magenta):

  \begin{itemize}
  \tightlist
  \item
    \(w^+(N_3) = \{(d, t), (e, t)\} \implies c(w^+(N_3)) = 5 + 5 = 10\).
  \end{itemize}
\end{enumerate}
\end{frame}

\begin{frame}{Ejemplo: Criterio de optimalidad}
\phantomsection\label{ejemplo-criterio-de-optimalidad}
\begin{itemize}
\tightlist
\item
  \textbf{Comparación entre flujo y capacidad de corte}:

  \begin{itemize}
  \tightlist
  \item
    El flujo de \(s\) a \(t\) transita con \(f_0 = 8\).
  \item
    La capacidad del Corte 1 es \(c(w^+(N_1)) = 8\).
  \end{itemize}
\item
  \textbf{Conclusión de optimalidad}:

  \begin{itemize}
  \tightlist
  \item
    Como \(f_0 = 8 = c(w^+(N_1))\), por la Proposición de Acotación:

    \begin{itemize}
    \tightlist
    \item
      El flujo \(\mathbf{f}\) es un \textbf{flujo máximo global}.
    \item
      El corte \(w^+(N_1)\) es un \textbf{corte de capacidad mínima}.
    \end{itemize}
  \end{itemize}
\end{itemize}
\end{frame}

\begin{frame}{Corolario: Criterio de optimalidad simultánea}
\phantomsection\label{corolario-criterio-de-optimalidad-simultuxe1nea}
\begin{tcolorbox}[enhanced jigsaw, coltitle=black, toprule=.15mm, opacitybacktitle=0.6, toptitle=1mm, leftrule=.75mm, breakable, left=2mm, opacityback=0, colback=white, bottomtitle=1mm, titlerule=0mm, title=\textcolor{quarto-callout-note-color}{\faInfo}\hspace{0.5em}{Criterio de optimalidad de flujo máximo y corte mínimo}, colbacktitle=quarto-callout-note-color!10!white, rightrule=.15mm, bottomrule=.15mm, arc=.35mm, colframe=quarto-callout-note-color-frame]

Si un flujo compatible \(\mathbf{f}\) y un \((s, t)\)-corte \(w^+(N)\)
satisfacen \(f_0 = c(w^+(N))\), entonces \(\mathbf{f}\) es el
\textbf{flujo máximo} de \(s\) a \(t\) y \(w^+(N)\) es el \textbf{corte
de mínima capacidad}.

\end{tcolorbox}

\begin{itemize}
\tightlist
\item
  Ningún otro flujo puede superar \(c(w^+(N)) = f_0\), y ningún otro
  corte puede tener capacidad inferior a \(f_0\).
\end{itemize}
\end{frame}

\begin{frame}{Corolario: Existencia y acotación de solución}
\phantomsection\label{corolario-existencia-y-acotaciuxf3n-de-soluciuxf3n}
\begin{tcolorbox}[enhanced jigsaw, coltitle=black, toprule=.15mm, opacitybacktitle=0.6, toptitle=1mm, leftrule=.75mm, breakable, left=2mm, opacityback=0, colback=white, bottomtitle=1mm, titlerule=0mm, title=\textcolor{quarto-callout-note-color}{\faInfo}\hspace{0.5em}{Existencia de solución finita}, colbacktitle=quarto-callout-note-color!10!white, rightrule=.15mm, bottomrule=.15mm, arc=.35mm, colframe=quarto-callout-note-color-frame]

El problema del flujo máximo de \(s\) a \(t\) tiene solución finita si y
sólo si no existe un camino orientado de \(s\) a \(t\) formado
íntegramente por arcos de capacidad infinita
(\(c_u = \infty, \forall u \in \mu\)).

\end{tcolorbox}

\begin{itemize}
\tightlist
\item
  Si no existe tal camino, el conjunto \(N\) de alcanzables desde \(s\)
  en \(G'=(V, U_\infty)\) forma un corte \(w^+(N)\) con
  \(c(w^+(N)) < \infty \implies f_0 \le c(w^+(N)) < \infty\).
\end{itemize}
\end{frame}

\begin{frame}{Teorema: Flujo máximo - corte mínimo}
\phantomsection\label{teorema-flujo-muxe1ximo---corte-muxednimo}
\begin{tcolorbox}[enhanced jigsaw, coltitle=black, toprule=.15mm, opacitybacktitle=0.6, toptitle=1mm, leftrule=.75mm, breakable, left=2mm, opacityback=0, colback=white, bottomtitle=1mm, titlerule=0mm, title=\textcolor{quarto-callout-important-color}{\faExclamation}\hspace{0.5em}{Flujo máximo - corte mínimo (Ford-Fulkerson, 1956)}, colbacktitle=quarto-callout-important-color!10!white, rightrule=.15mm, bottomrule=.15mm, arc=.35mm, colframe=quarto-callout-important-color-frame]

Dada una red de transporte \((G = (V, U), \mathbf{c})\) y dos vértices
\(s, t \in V\), el \((s, t)\)-flujo máximo coincide exactamente con la
mínima capacidad de los \((s, t)\)-cortes:

\[
\max_{\mathbf{f}} f_0 = \min_{\substack{N \subset V \\ s \in N, \, t \notin N}} c(w^+(N))
\]

\end{tcolorbox}
\end{frame}

\begin{frame}{Interpretación combinatoria de Ford-Fulkerson}
\phantomsection\label{interpretaciuxf3n-combinatoria-de-ford-fulkerson}
\begin{itemize}
\tightlist
\item
  \textbf{Dualidad Min-Max en redes}:

  \begin{itemize}
  \tightlist
  \item
    Establece la igualdad estricta entre la máxima cantidad de flujo
    transmisible y la mínima capacidad que puede bloquear la red.
  \item
    Garantiza la existencia de una barrera saturada cuando el flujo
    alcanza su valor óptimo global.
  \end{itemize}
\end{itemize}
\end{frame}

\begin{frame}{Ejemplo Minty: Planteamiento y Coloración de Arcos}
\phantomsection\label{ejemplo-minty-planteamiento-y-coloraciuxf3n-de-arcos}
\begin{itemize}
\tightlist
\item
  Considérese la red de 6 nodos \(V = \{s, 1, 2, 3, 4, t\}\) con flujo
  \(f_0 = 10\) y retorno \(u_0 = (t, s)\):
\item
  \textbf{Reglas de coloración según caudales actuales}:

  \begin{itemize}
  \tightlist
  \item
    Arco de retorno \(u_0 = (t, s)\) (\(f_0 = 10, c_0 = \infty\)):
    \textbf{Negro} (admite aumento).
  \item
    Arco \((3, 1)\) (\(f_{31} = 0, c_{31} = 5\)): \textbf{Negro}
    (\(0 = f_{31} < c_{31}\)).
  \item
    Arcos a capacidad máxima: \((1,2)\), \((1,4)\), \((3,4)\),
    \((4,2)\): \textbf{Verdes} (\(0 < f_u = c_u\)).
  \item
    Arcos a capacidad intermedia: \((s,1)\), \((s,3)\), \((2,t)\),
    \((4,t)\): \textbf{Rojos} (\(0 < f_u < c_u\)).
  \end{itemize}
\end{itemize}
\end{frame}

\begin{frame}{Ejemplo Minty: Gráfico y Marcado de Vértices}
\phantomsection\label{ejemplo-minty-gruxe1fico-y-marcado-de-vuxe9rtices}
\begin{itemize}
\tightlist
\item
  \textbf{Propagación del marcado desde la fuente}:

  \begin{itemize}
  \tightlist
  \item
    Inicio: \(S = \{s\}\).
  \item
    Vía arcos rojos salientes \((s,1)\) y
    \((s,3) \implies S = \{s, 1, 3\}\).
  \item
    Arcos salientes de \(1\) y \(3\) son verdes (bloqueados); entrantes
    son negros \(\implies\) Marcado final \(S = \{s, 1, 3\}\).
  \end{itemize}
\end{itemize}
\end{frame}

\begin{frame}{Ejemplo Minty: Verificación de Optimalidad}
\phantomsection\label{ejemplo-minty-verificaciuxf3n-de-optimalidad}
\begin{itemize}
\tightlist
\item
  \textbf{Análisis de la partición \(S = \{s, 1, 3\}\)}:

  \begin{itemize}
  \tightlist
  \item
    El sumidero \(t \notin S \implies\) No existe cadena incremental de
    \(s\) a \(t\).
  \item
    La partición \(N = \{s, 1, 3\}\) define un \((s, t)\)-corte válido.
  \end{itemize}
\item
  \textbf{Capacidad del corte \(w^+(N)\)}:

  \begin{itemize}
  \tightlist
  \item
    Arcos salientes \(w^+(N) = \{(1,2), (1,4), (3,4)\}\) (todos verdes).
  \item
    \(c(w^+(N)) = c_{12} + c_{14} + c_{34} = 4 + 1 + 5 = 10\).
  \end{itemize}
\item
  \textbf{Conclusión}: \(f_0 = 10 = c(w^+(N)) \implies\) El flujo
  \(\mathbf{f}\) es \textbf{óptimo global} y \(w^+(N)\) es \textbf{corte
  de capacidad mínima}.
\end{itemize}
\end{frame}

\begin{frame}{Ejemplo Minty: Análisis de Sensibilidad (Red62)}
\phantomsection\label{ejemplo-minty-anuxe1lisis-de-sensibilidad-red62}
\begin{itemize}
\tightlist
\item
  \textbf{Modificación de capacidad}: Al elevar \(c_{14} = 6\)
  manteniéndose \(f_{14} = 1\):

  \begin{itemize}
  \tightlist
  \item
    El arco \((1,4)\) pasa de verde a \textbf{rojo} (\(0 < 1 < 6\)).
  \item
    Desde \(1 \in S\), el arco rojo \((1,4)\) marca el nodo \(4\), y el
    rojo \((4,t)\) marca finalmente el sumidero \(t \implies\) ¡Se
    habilita el incremento de flujo vía \(s \to 1 \to 4 \to t\)!
  \end{itemize}
\end{itemize}
\end{frame}

\begin{frame}{Algoritmo de Ford-Fulkerson: Introducción}
\phantomsection\label{algoritmo-de-ford-fulkerson-introducciuxf3n}
\begin{itemize}
\tightlist
\item
  \textbf{Orígenes y fundamento}:

  \begin{itemize}
  \tightlist
  \item
    Publicado en 1956 por L. R. Ford (1927--2017) y D. R. Fulkerson
    (1924--1976).
  \item
    Denominado también \textbf{algoritmo del camino de aumento}.
  \end{itemize}
\item
  \textbf{Idea central}:

  \begin{itemize}
  \tightlist
  \item
    Incrementar progresivamente el caudal de la red a lo largo de
    trayectorias orientadas no saturadas (cadenas incrementales) desde
    la fuente \(s\) hasta el sumidero \(t\).
  \end{itemize}
\item
  \textbf{Criterio de parada y optimalidad}:

  \begin{itemize}
  \tightlist
  \item
    El algoritmo finaliza cuando no es posible hallar ninguna
    trayectoria adicional de aumento.
  \end{itemize}
\end{itemize}
\end{frame}

\begin{frame}{Definición: Red residual (Grafo incremental)}
\phantomsection\label{definiciuxf3n-red-residual-grafo-incremental}
\begin{tcolorbox}[enhanced jigsaw, coltitle=black, toprule=.15mm, opacitybacktitle=0.6, toptitle=1mm, leftrule=.75mm, breakable, left=2mm, opacityback=0, colback=white, bottomtitle=1mm, titlerule=0mm, title=\textcolor{quarto-callout-note-color}{\faInfo}\hspace{0.5em}{Red residual G(f)}, colbacktitle=quarto-callout-note-color!10!white, rightrule=.15mm, bottomrule=.15mm, arc=.35mm, colframe=quarto-callout-note-color-frame]

Dado un flujo compatible \(\mathbf{f}\) en la red de transporte
\((G = (V, U), \mathbf{c})\), se define el \textbf{grafo incremental} (o
\textbf{red residual}) \(G(\mathbf{f}) = (V, U(\mathbf{f}))\) como el
digrafo donde \(U(\mathbf{f})\) contiene:

\[
\begin{cases}
u^+ = (i, j) \in U(\mathbf{f}) \text{ con peso } \ell_{u^+} = c_u - f_u, & \text{si } f_u < c_u \\
u^- = (j, i) \in U(\mathbf{f}) \text{ con peso } \ell_{u^-} = f_u, & \text{si } f_u > 0
\end{cases}
\]

\end{tcolorbox}
\end{frame}

\begin{frame}{Propiedades del grafo incremental G(f)}
\phantomsection\label{propiedades-del-grafo-incremental-gf}
\begin{itemize}
\tightlist
\item
  \textbf{Doble arco simultáneo}:

  \begin{itemize}
  \tightlist
  \item
    Si un arco satisface \(0 < f_u < c_u\), genera tanto un arco directo
    \(u^+\) como un arco inverso \(u^-\) en \(G(\mathbf{f})\).
  \end{itemize}
\item
  \textbf{Arco ficticio de retorno}:

  \begin{itemize}
  \tightlist
  \item
    El arco de retorno \(u_0 = (t, s)\) \textbf{no} forma parte del
    grafo incremental.
  \end{itemize}
\item
  \textbf{Teorema de caracterización}:

  \begin{itemize}
  \tightlist
  \item
    Un flujo compatible \(\mathbf{f}\) es \textbf{óptimo} \(\iff\)
    \textbf{no existe} ningún camino orientado de \(s\) a \(t\) en la
    red residual \(G(\mathbf{f})\).
  \end{itemize}
\end{itemize}
\end{frame}

\begin{frame}{Ejemplo: Construcción del Grafo Incremental (Reglas)}
\phantomsection\label{ejemplo-construcciuxf3n-del-grafo-incremental-reglas}
\begin{itemize}
\item
  Red de 6 nodos \(V = \{s, 2, 3, 4, 5, t\}\) con flujo de valor
  \(f_0 = 3\):
\item
  \textbf{Arcos no saturados (\(f_u < c_u\))}:

  \begin{itemize}
  \tightlist
  \item
    Para el arco \((s,2)\) con \(f=1, c=3 \implies\) genera arco directo
    \((s,2)\) con peso residual \(\ell = 3 - 1 = 2\).
  \end{itemize}
\item
  \textbf{Arcos con flujo positivo (\(f_u > 0\))}:

  \begin{itemize}
  \tightlist
  \item
    Para el mismo arco \((s,2)\) con \(f=1 \implies\) genera arco
    inverso \((2,s)\) con peso residual \(\ell = 1\).
  \end{itemize}
\end{itemize}
\end{frame}

\begin{frame}{Ejemplo: Grafo Incremental (Saturados y Nulos)}
\phantomsection\label{ejemplo-grafo-incremental-saturados-y-nulos}
\begin{itemize}
\tightlist
\item
  \textbf{Arcos a capacidad máxima (\(f_u = c_u\))}:

  \begin{itemize}
  \tightlist
  \item
    Para el arco \((s,3)\) con \(f=2, c=2 \implies\) no admite
    incremento (\(u^+\) no existe). Genera únicamente el arco inverso
    \((3,s)\) con peso \(\ell = 2\).
  \end{itemize}
\item
  \textbf{Arcos con flujo nulo (\(f_u = 0\))}:

  \begin{itemize}
  \tightlist
  \item
    Para el arco \((3,2)\) con \(f=0, c=1 \implies\) no admite reducción
    (\(u^-\) no existe). Genera únicamente el arco directo \((3,2)\) con
    peso \(\ell = 1\).
  \end{itemize}
\end{itemize}
\end{frame}

\begin{frame}{Ejemplo: Gráfico de la Red Residual}
\phantomsection\label{ejemplo-gruxe1fico-de-la-red-residual}
\begin{itemize}
\tightlist
\item
  Transformación de una red de transporte con flujo \(\mathbf{f}\)
  (izquierda) a su grafo incremental \(G(\mathbf{f})\) con pesos
  residuales \(\ell_u\) (derecha).
\end{itemize}
\end{frame}

\begin{frame}{Esquema: Algoritmo por Red Residual (Pasos 1 y 2)}
\phantomsection\label{esquema-algoritmo-por-red-residual-pasos-1-y-2}
\begin{itemize}
\tightlist
\item
  \textbf{Paso 1 (Inicialización)}:

  \begin{itemize}
  \tightlist
  \item
    Hacer \(k = 0\). Determinar un flujo inicial compatible
    \(\mathbf{f}^0 = \mathbf{0}\) (\(f_0^0 = 0\)).
  \end{itemize}
\item
  \textbf{Paso 2 (Construcción y Búsqueda)}:

  \begin{itemize}
  \tightlist
  \item
    Construir el grafo incremental \(G(\mathbf{f}^k)\). Buscar un camino
    orientado \(\mu^k\) desde \(s\) hasta \(t\) en \(G(\mathbf{f}^k)\).
  \item
    \textbf{Si no existe camino}: \textbf{PARAR}. El flujo actual
    \(\mathbf{f}^k\) es \textbf{óptimo} (y el conjunto de alcanzables
    desde \(s\) define el \textbf{corte mínimo}).
  \item
    \textbf{Si existe camino}: Ir al Paso 3.
  \end{itemize}
\end{itemize}
\end{frame}

\begin{frame}{Esquema: Algoritmo por Red Residual (Paso 3)}
\phantomsection\label{esquema-algoritmo-por-red-residual-paso-3}
\begin{itemize}
\tightlist
\item
  \textbf{Paso 3 (Aumento de Flujo y Actualización)}:

  \begin{itemize}
  \tightlist
  \item
    Cuello de botella en la cadena incremental:
    \[\epsilon^k = \min \{\ell_u \mid u \in \mu^k\}\]
  \item
    Actualización de flujos sobre cada arco \(u \in U\):
    \[f_u^{k+1} = \begin{cases} f_u^k + \epsilon^k, & \text{si } u \in \mu^k(+) \text{ (directo en la red)} \\ f_u^k - \epsilon^k, & \text{si } u \in \mu^k(-) \text{ (inverso en la red)} \\ f_u^k, & \text{si } u \notin \mu^k \end{cases}\]
  \item
    Incrementar flujo de retorno \(f_0^{k+1} = f_0^k + \epsilon^k\),
    hacer \(k \leftarrow k+1\) y volver al Paso 2.
  \end{itemize}
\end{itemize}
\end{frame}

\begin{frame}{Ejemplo Ford-Fulkerson: Planteamiento y Red Inicial}
\phantomsection\label{ejemplo-ford-fulkerson-planteamiento-y-red-inicial}
\begin{itemize}
\tightlist
\item
  \textbf{Red de 6 nodos} \(V = \{s, 2, 3, 4, 5, t\}\) con fuente \(s\),
  sumidero \(t\) y capacidades:

  \begin{itemize}
  \tightlist
  \item
    \(c_{s2}=3, \quad c_{s3}=2, \quad c_{32}=1, \quad c_{24}=2\)
  \item
    \(c_{43}=1, \quad c_{35}=2, \quad c_{54}=1, \quad c_{4t}=3, \quad c_{5t}=2\)
  \end{itemize}
\item
  \textbf{Estado Inicial (\(k=0\))}:

  \begin{itemize}
  \tightlist
  \item
    Partimos del flujo trivial nulo \(\mathbf{f}^0 = \mathbf{0}\)
    (\(f_0^0 = 0\)). Todos los arcos se encuentran a flujo 0.
  \end{itemize}
\end{itemize}
\end{frame}

\begin{frame}{Ejemplo Ford-Fulkerson: Iteración 1 (Búsqueda)}
\phantomsection\label{ejemplo-ford-fulkerson-iteraciuxf3n-1-buxfasqueda}
\begin{itemize}
\tightlist
\item
  \textbf{Construcción de la Red Residual \(G(\mathbf{f}^0)\)}:

  \begin{itemize}
  \tightlist
  \item
    Al ser \(\mathbf{f}^0 = \mathbf{0}\), \(G(\mathbf{f}^0)\) contiene
    únicamente los arcos directos \(u^+\) con capacidades residuales
    \(\ell_{u^+} = c_u\).
  \end{itemize}
\item
  \textbf{Búsqueda de Camino Aumentante}:

  \begin{itemize}
  \tightlist
  \item
    Se selecciona la trayectoria orientada
    \(\mu_1 = (s \to 2 \to 4 \to t)\).
  \end{itemize}
\item
  \textbf{Cálculo del Cuello de Botella}:
  \[\epsilon_1 = \min \{\ell_{s2}, \ell_{24}, \ell_{4t}\} = \min \{3, 2, 3\} = 2\]
\end{itemize}
\end{frame}

\begin{frame}{Ejemplo Ford-Fulkerson: Iteración 1 (Actualización)}
\phantomsection\label{ejemplo-ford-fulkerson-iteraciuxf3n-1-actualizaciuxf3n}
\begin{itemize}
\tightlist
\item
  \textbf{Actualización del vector de flujo \(\mathbf{f}^1\)}:

  \begin{itemize}
  \tightlist
  \item
    \(f_{s2}^1 = 0 + 2 = 2\)
  \item
    \(f_{24}^1 = 0 + 2 = 2 \implies\) Arco \((2,4)\) \textbf{saturado}
    (pasa a verde)
  \item
    \(f_{4t}^1 = 0 + 2 = 2\)
  \end{itemize}
\item
  \textbf{Caudal de retorno y estado de la red}:

  \begin{itemize}
  \tightlist
  \item
    Flujo circulante de retorno: \(f_0^1 = 0 + 2 = 2\).
  \end{itemize}
\end{itemize}
\end{frame}

\begin{frame}{Ejemplo Ford-Fulkerson: Iteración 2 (Búsqueda)}
\phantomsection\label{ejemplo-ford-fulkerson-iteraciuxf3n-2-buxfasqueda}
\begin{itemize}
\tightlist
\item
  \textbf{Construcción de la Red Residual \(G(\mathbf{f}^1)\)}:

  \begin{itemize}
  \tightlist
  \item
    Al estar saturado el arco \((2,4)\), solo genera en
    \(G(\mathbf{f}^1)\) el arco inverso \((4,2)\) con peso \(\ell = 2\).
  \item
    Arco \((s,2)\) admite incremento \(\ell = 3 - 2 = 1\) e inverso
    \((2,s)\) con peso \(\ell = 2\).
  \end{itemize}
\item
  \textbf{Búsqueda de Camino Aumentante}:

  \begin{itemize}
  \tightlist
  \item
    Se selecciona la trayectoria orientada
    \(\mu_2 = (s \to 3 \to 5 \to t)\).
  \end{itemize}
\item
  \textbf{Cálculo del Cuello de Botella}:
  \[\epsilon_2 = \min \{\ell_{s3}, \ell_{35}, \ell_{5t}\} = \min \{2, 2, 2\} = 2\]
\end{itemize}
\end{frame}

\begin{frame}{Ejemplo Ford-Fulkerson: Iteración 2 (Actualización)}
\phantomsection\label{ejemplo-ford-fulkerson-iteraciuxf3n-2-actualizaciuxf3n}
\begin{itemize}
\tightlist
\item
  \textbf{Actualización del vector de flujo \(\mathbf{f}^2\)}:

  \begin{itemize}
  \tightlist
  \item
    \(f_{s3}^2 = 0 + 2 = 2 \implies\) Arco \((s,3)\) \textbf{saturado}
    (verde)
  \item
    \(f_{35}^2 = 0 + 2 = 2 \implies\) Arco \((3,5)\) \textbf{saturado}
    (verde)
  \item
    \(f_{5t}^2 = 0 + 2 = 2 \implies\) Arco \((5,t)\) \textbf{saturado}
    (verde)
  \end{itemize}
\item
  \textbf{Caudal total de retorno}:

  \begin{itemize}
  \tightlist
  \item
    Flujo circulante de retorno: \(f_0^2 = 2 + 2 = 4\).
  \end{itemize}
\end{itemize}
\end{frame}

\begin{frame}{Ejemplo Ford-Fulkerson: Iteración 3 (Criterio de Parada)}
\phantomsection\label{ejemplo-ford-fulkerson-iteraciuxf3n-3-criterio-de-parada}
\begin{itemize}
\tightlist
\item
  \textbf{Evaluación de la Red Residual \(G(\mathbf{f}^2)\)}:

  \begin{itemize}
  \tightlist
  \item
    Todos los arcos salientes desde la fuente \(s\) a través de \(3\)
    están saturados.
  \item
    El único nodo alcanzable desde \(s\) mediante caminos orientados en
    \(G(\mathbf{f}^2)\) es el nodo \(2\).
  \end{itemize}
\item
  \textbf{Conclusión de Parada}:

  \begin{itemize}
  \tightlist
  \item
    ¡No existe ningún camino orientado desde \(s\) hasta \(t\) en
    \(G(\mathbf{f}^2)\)!
  \item
    El flujo \(\mathbf{f}^2\) es el \textbf{flujo máximo global} con
    valor \(f_0^* = 4\).
  \end{itemize}
\end{itemize}
\end{frame}

\begin{frame}{Ejemplo Ford-Fulkerson: Corte de Capacidad Mínima}
\phantomsection\label{ejemplo-ford-fulkerson-corte-de-capacidad-muxednima}
\begin{itemize}
\item
  \textbf{Nodos alcanzables desde la fuente \(s\) en
  \(G(\mathbf{f}^2)\)}: \[N = \{s, 2\}\]
\item
  \textbf{Corte de capacidad mínima resultante}:
  \[w^+(N) = \{(s,3), (2,4)\}\]
\item
  \textbf{Capacidad del corte de mínima capacidad}:
  \[C(w^+(N)) = c_{s3} + c_{24} = 2 + 2 = 4 = f_0^*\]

  \begin{itemize}
  \tightlist
  \item
    Verifica de forma exacta el Teorema de Ford-Fulkerson
    (\(f_0^* = C(w^+(N))\)).
  \end{itemize}
\end{itemize}
\end{frame}

\begin{frame}{Ejemplo Ford-Fulkerson: Reproductor Interactivo}
\phantomsection\label{ejemplo-ford-fulkerson-reproductor-interactivo}
\end{frame}

\begin{frame}{Ejemplo Ford-Fulkerson: Código de Colores}
\phantomsection\label{ejemplo-ford-fulkerson-cuxf3digo-de-colores}
\begin{itemize}
\tightlist
\item
  \textbf{Significado visual de la simbología en los gráficos}:

  \begin{itemize}
  \tightlist
  \item
    \textbf{Arcos en naranja (\(\color{#f97316}{\rightarrow}\))}: Camino
    aumentante seleccionado en la iteración actual.
  \item
    \textbf{Arcos en verde (\(\color{#16a34a}{\rightarrow}\))}: Arcos
    saturados a capacidad máxima (\(f_u = c_u\)).
  \item
    \textbf{Arcos en rojo (\(\color{firebrick}{\rightarrow}\))}: Flujo
    parcial no nulo (\(0 < f_u < c_u\)).
  \item
    \textbf{Arcos en negro (\(\rightarrow\))}: Flujo nulo (\(f_u = 0\)).
  \item
    \textbf{Arcos/Nodos en azul (\(\color{#2563eb}{\rightarrow}\))}:
    Destacan los nodos alcanzables \(N = \{s, 2\}\) y el \textbf{corte
    mínimo} \(w^+(N)\).
  \end{itemize}
\end{itemize}
\end{frame}

\begin{frame}{Algoritmo por Etiquetado de Vértices: Motivación}
\phantomsection\label{algoritmo-por-etiquetado-de-vuxe9rtices-motivaciuxf3n}
\begin{itemize}
\tightlist
\item
  \textbf{Operación directa sobre la red original}:

  \begin{itemize}
  \tightlist
  \item
    Evita la necesidad de construir e ilustrar explícitamente la red
    residual \(G(\mathbf{f})\) en cada iteración.
  \item
    Emplea una \textbf{regla de etiquetado dinámico} que opera asignando
    marcas directamente sobre los vértices de la red.
  \end{itemize}
\item
  \textbf{Marca de la fuente}:

  \begin{itemize}
  \tightlist
  \item
    Se asigna inicialmente
    \(\text{Etiqueta}(s) = (\text{–}, \, \color{#0284c7}{\delta_s = +\infty})\)
    y se propaga la información de cuello de botella a los nodos vecinos
    no explorados.
  \end{itemize}
\end{itemize}
\end{frame}

\begin{frame}{Definición: Estructura de las etiquetas de vértices}
\phantomsection\label{definiciuxf3n-estructura-de-las-etiquetas-de-vuxe9rtices}
\begin{tcolorbox}[enhanced jigsaw, coltitle=black, toprule=.15mm, opacitybacktitle=0.6, toptitle=1mm, leftrule=.75mm, breakable, left=2mm, opacityback=0, colback=white, bottomtitle=1mm, titlerule=0mm, title=\textcolor{quarto-callout-note-color}{\faInfo}\hspace{0.5em}{Estructura de las etiquetas de vértices}, colbacktitle=quarto-callout-note-color!10!white, rightrule=.15mm, bottomrule=.15mm, arc=.35mm, colframe=quarto-callout-note-color-frame]

La etiqueta asignada a un nodo \(j\) presenta dos componentes
diferenciadas:

\[
\text{Etiqueta}(j) = (\color{#0f172a}{\pm i}, \, \color{#0284c7}{\delta_j})
\]

\begin{enumerate}
\tightlist
\item
  \textbf{Procedencia (\(\color{#0f172a}{\pm i}\))}: Nodo antecedente
  \(i\) (ya etiquetado). Marca \textbf{\(\color{#16a34a}{+i}\)} indica
  arco directo saliente \((i, j)\) con \(f_{ij} < c_{ij}\); marca
  \textbf{\(\color{firebrick}{-i}\)} indica arco inverso entrante
  \((j, i)\) con \(f_{ji} > 0\).
\item
  \textbf{Capacidad disponible (\(\color{#0284c7}{\delta_j}\))}: Caudal
  máximo incremental transmisible desde \(s\) hasta \(j\).
\end{enumerate}

\end{tcolorbox}
\end{frame}

\begin{frame}{Definición: Reglas formales de etiquetado}
\phantomsection\label{definiciuxf3n-reglas-formales-de-etiquetado}
\begin{tcolorbox}[enhanced jigsaw, coltitle=black, toprule=.15mm, opacitybacktitle=0.6, toptitle=1mm, leftrule=.75mm, breakable, left=2mm, opacityback=0, colback=white, bottomtitle=1mm, titlerule=0mm, title=\textcolor{quarto-callout-note-color}{\faInfo}\hspace{0.5em}{Reglas formales de etiquetado}, colbacktitle=quarto-callout-note-color!10!white, rightrule=.15mm, bottomrule=.15mm, arc=.35mm, colframe=quarto-callout-note-color-frame]

Sea \(i \in V\) etiquetado con \((\cdot, \, \color{#0284c7}{\delta_i})\)
y sea \(j \in V\) adyacente aún \textbf{no etiquetado}:

\begin{enumerate}
\item
  \textbf{Regla 1 (Arco Directo \((i, j)\))}: Si \(f_{ij} < c_{ij}\),
  etiquetar \(j\) como
  \((\color{#16a34a}{+i}, \, \color{#0284c7}{\delta_j})\) con:
  \[ \color{#0284c7}{\delta_j} = \min \left\{ \color{#0284c7}{\delta_i}, \, c_{ij} - f_{ij} \right\} \]
\item
  \textbf{Regla 2 (Arco Inverso \((j, i)\))}: Si \(f_{ji} > 0\),
  etiquetar \(j\) como
  \((\color{firebrick}{-i}, \, \color{#0284c7}{\delta_j})\) con:
  \[ \color{#0284c7}{\delta_j} = \min \left\{ \color{#0284c7}{\delta_i}, \, f_{ji} \right\} \]
\end{enumerate}

\end{tcolorbox}
\end{frame}

\begin{frame}{Ilustración: Reglas formales de etiquetado}
\phantomsection\label{ilustraciuxf3n-reglas-formales-de-etiquetado}
\begin{itemize}
\tightlist
\item
  Asignación de marcas directas \((+i, \delta_j)\) en arcos salientes
  (izquierda) y marcas inversas \((-i, \delta_j)\) en arcos entrantes
  (derecha).
\end{itemize}
\end{frame}

\begin{frame}{Esquema: Algoritmo por Etiquetado (Pasos 1 a 4)}
\phantomsection\label{esquema-algoritmo-por-etiquetado-pasos-1-a-4}
\begin{itemize}
\tightlist
\item
  \textbf{Paso 1 (Inicialización)}: Asignar flujo inicial
  \(\mathbf{f} = \mathbf{0}\). Etiquetar fuente \(s\) como
  \((\text{–}, \, \color{#0284c7}{+\infty})\).
\item
  \textbf{Paso 2 (Evaluación)}: Comprobar si todos los nodos etiquetados
  han sido explorados.
\item
  \textbf{Paso 3 (Propagación)}: Etiquetar nodos adyacentes posibles
  mediante la \textbf{Regla 1} (directos) y \textbf{Regla 2} (inversos).
\item
  \textbf{Paso 4 (Verificación del Sumidero)}:

  \begin{itemize}
  \tightlist
  \item
    Si \(t\) resulta etiquetado \(\implies\) Fijar la cantidad de
    aumento \(\Delta = \color{#0284c7}{\delta_t}\) e ir al Paso 5.
  \item
    Si \(t\) no resulta etiquetado \(\implies\) Ir al Paso 6.
  \end{itemize}
\end{itemize}
\end{frame}

\begin{frame}{Esquema: Algoritmo por Etiquetado (Pasos 5 y 6)}
\phantomsection\label{esquema-algoritmo-por-etiquetado-pasos-5-y-6}
\begin{itemize}
\tightlist
\item
  \textbf{Paso 5 (Modificación del Flujo)}:

  \begin{itemize}
  \tightlist
  \item
    Rastrear desde \(t\) hacia atrás mediante los antecedentes:

    \begin{itemize}
    \tightlist
    \item
      Si la marca es \(+i \implies f_{ik} \leftarrow f_{ik} + \Delta\).
    \item
      Si la marca es \(-i \implies f_{ki} \leftarrow f_{ki} - \Delta\).
    \end{itemize}
  \item
    Al llegar a \(s\): hacer \(f_0 \leftarrow f_0 + \Delta\), borrar
    marcas de los nodos (salvo \(s\)) y regresar al Paso 2.
  \end{itemize}
\item
  \textbf{Paso 6 (Optimalidad y Parada)}:

  \begin{itemize}
  \tightlist
  \item
    \textbf{PARAR}. El flujo actual es \textbf{óptimo} (\(f_0^*\)). El
    conjunto de nodos etiquetados \(N = S\) define el \textbf{corte de
    capacidad mínima} \(w^+(N)\).
  \end{itemize}
\end{itemize}
\end{frame}

\begin{frame}{Relación entre Etiquetado de Vértices y BFS}
\phantomsection\label{relaciuxf3n-entre-etiquetado-de-vuxe9rtices-y-bfs}
\begin{itemize}
\tightlist
\item
  \textbf{Búsqueda en Anchura (BFS)}:

  \begin{itemize}
  \tightlist
  \item
    La propagación de marcas por niveles explora la red en anchura desde
    la fuente \(s\).
  \item
    Garantiza que el camino de aumento seleccionado en cada iteración
    sea el de \textbf{menor longitud en número de arcos}.
  \end{itemize}
\item
  \textbf{Ventaja Algorítmica}:

  \begin{itemize}
  \tightlist
  \item
    Evita la exploración de trayectorias largas e ineficientes que
    ralentizan la convergencia.
  \end{itemize}
\end{itemize}
\end{frame}

\begin{frame}{Ejemplo Etiquetado: Planteamiento y Capacidades}
\phantomsection\label{ejemplo-etiquetado-planteamiento-y-capacidades}
\begin{itemize}
\tightlist
\item
  \textbf{Red de 8 nodos} \(V = \{s, 1, 2, 3, 4, 5, 6, t\}\) con
  capacidades entre paréntesis:

  \begin{itemize}
  \tightlist
  \item
    \(c_{s1} = 7\), \(c_{s2} = 2\), \(c_{12} = 1\), \(c_{13} = 5\),
    \(c_{32} = 3\), \(c_{34} = 2\)
  \item
    \(c_{35} = 4\), \(c_{24} = 4\), \(c_{46} = 5\), \(c_{56} = 3\),
    \(c_{5t} = 3\), \(c_{6t} = 6\)
  \end{itemize}
\item
  \textbf{Estado Inicial}: Flujo nulo \(\mathbf{f}^0 = \mathbf{0}\)
  (\(f_0^0 = 0\)).
\end{itemize}
\end{frame}

\begin{frame}{Ejemplo Etiquetado: Iteración 1 (Marcado de Nodos)}
\phantomsection\label{ejemplo-etiquetado-iteraciuxf3n-1-marcado-de-nodos}
\begin{itemize}
\tightlist
\item
  \textbf{Propagación de marcas desde la fuente \(s\)}:

  \begin{itemize}
  \tightlist
  \item
    Etiqueta inicial:
    \(\text{Etiqueta}(s) = (\text{–}, \, \color{#0284c7}{\delta_s = +\infty})\).
  \item
    Desde \(s\): etiquetar \(1\) con
    \((\color{#16a34a}{+s}, \, \color{#0284c7}{7})\) y \(2\) con
    \((\color{#16a34a}{+s}, \, \color{#0284c7}{2})\).
  \item
    Desde \(2\): etiquetar \(4\) con
    \((\color{#16a34a}{+2}, \, \color{#0284c7}{\min\{2, 4\} = 2})\).
  \item
    Desde \(4\): etiquetar \(6\) con
    \((\color{#16a34a}{+4}, \, \color{#0284c7}{\min\{2, 5\} = 2})\).
  \item
    Desde \(6\): etiquetar sumidero \(t\) con
    \((\color{#16a34a}{+6}, \, \color{#0284c7}{\min\{2, 6\} = 2})\).
  \end{itemize}
\end{itemize}
\end{frame}

\begin{frame}{Ejemplo Etiquetado: Iteración 1 (Aumento y Actualización)}
\phantomsection\label{ejemplo-etiquetado-iteraciuxf3n-1-aumento-y-actualizaciuxf3n}
\begin{itemize}
\item
  \textbf{Camino de aumento detectado}:
  \[\mu_1 = (s \to 2 \to 4 \to 6 \to t) \quad \text{con } \Delta_1 = \color{#0284c7}{\delta_t = 2}\]
\item
  \textbf{Actualización de caudales (\(\mathbf{f}^1\))}:

  \begin{itemize}
  \tightlist
  \item
    \(f_{s2} \leftarrow 0 + 2 = 2 \implies\) Arco \((s,2)\)
    \textbf{saturado} (verde).
  \item
    \(f_{24} \leftarrow 0 + 2 = 2\), \(f_{46} \leftarrow 0 + 2 = 2\),
    \(f_{6t} \leftarrow 0 + 2 = 2\).
  \item
    Flujo de retorno circulante: \(f_0^1 = 2\).
  \end{itemize}
\end{itemize}
\end{frame}

\begin{frame}{Ejemplo Etiquetado: Iteración 2 (Marcado de Nodos)}
\phantomsection\label{ejemplo-etiquetado-iteraciuxf3n-2-marcado-de-nodos}
\begin{itemize}
\tightlist
\item
  \textbf{Propagación de marcas (Arco \((s,2)\) saturado, no permite
  etiquetar \(2\))}:

  \begin{itemize}
  \tightlist
  \item
    Etiqueta en \(s\): \((\text{–}, \, \color{#0284c7}{+\infty})\).
  \item
    Desde \(s\): etiquetar \(1\) con
    \((\color{#16a34a}{+s}, \, \color{#0284c7}{7})\).
  \item
    Desde \(1\): etiquetar \(3\) con
    \((\color{#16a34a}{+1}, \, \color{#0284c7}{\min\{7, 5\} = 5})\) y
    \(2\) con
    \((\color{#16a34a}{+1}, \, \color{#0284c7}{\min\{7, 1\} = 1})\).
  \item
    Desde \(3\): etiquetar \(5\) con
    \((\color{#16a34a}{+3}, \, \color{#0284c7}{\min\{5, 4\} = 4})\) y
    \(4\) con
    \((\color{#16a34a}{+3}, \, \color{#0284c7}{\min\{5, 2\} = 2})\).
  \item
    Desde \(5\): etiquetar sumidero \(t\) con
    \((\color{#16a34a}{+5}, \, \color{#0284c7}{\min\{4, 3\} = 3})\).
  \end{itemize}
\end{itemize}
\end{frame}

\begin{frame}{Ejemplo Etiquetado: Iteración 2 (Aumento y Actualización)}
\phantomsection\label{ejemplo-etiquetado-iteraciuxf3n-2-aumento-y-actualizaciuxf3n}
\begin{itemize}
\item
  \textbf{Camino de aumento detectado}:
  \[\mu_2 = (s \to 1 \to 3 \to 5 \to t) \quad \text{con } \Delta_2 = \color{#0284c7}{\delta_t = 3}\]
\item
  \textbf{Actualización de caudales (\(\mathbf{f}^2\))}:

  \begin{itemize}
  \tightlist
  \item
    \(f_{s1} \leftarrow 0 + 3 = 3\), \(f_{13} \leftarrow 0 + 3 = 3\),
    \(f_{35} \leftarrow 0 + 3 = 3\).
  \item
    \(f_{5t} \leftarrow 0 + 3 = 3 \implies\) Arco \((5,t)\)
    \textbf{saturado} (verde).
  \item
    Flujo de retorno circulante: \(f_0^2 = 2 + 3 = 5\).
  \end{itemize}
\end{itemize}
\end{frame}

\begin{frame}{Ejemplo Etiquetado: Iteraciones 3 y 4 (Resumen)}
\phantomsection\label{ejemplo-etiquetado-iteraciones-3-y-4-resumen}
\begin{itemize}
\tightlist
\item
  \textbf{Iteración 3}:

  \begin{itemize}
  \tightlist
  \item
    Camino de aumento: \(s \to 1 \to 3 \to 4 \to 6 \to t\) con
    \(\Delta_3 = 2\).
  \item
    \(f_0^3 = 5 + 2 = 7\). Arcos \((3,4)\) y \((1,3)\) quedan saturados.
  \end{itemize}
\item
  \textbf{Iteración 4}:

  \begin{itemize}
  \tightlist
  \item
    Búsqueda por vías alternativas y actualización de caudales
    residuales.
  \item
    Flujo circulante total acumulado: \(f_0^4 = 7\).
  \end{itemize}
\end{itemize}
\end{frame}

\begin{frame}{Ejemplo Etiquetado: Iteración 5 (Marcado de Nodos)}
\phantomsection\label{ejemplo-etiquetado-iteraciuxf3n-5-marcado-de-nodos}
\begin{itemize}
\tightlist
\item
  \textbf{Propagación de marcas
  (\(f_{s1} = 5, f_{s2} = 2, f_{13} = 5\))}:

  \begin{itemize}
  \tightlist
  \item
    Etiqueta en \(s\): \((\text{–}, \, \color{#0284c7}{+\infty})\).
  \item
    Desde \(s\): etiquetar \(1\) con
    \((\color{#16a34a}{+s}, \, \color{#0284c7}{7 - 5 = 2})\).
  \item
    Desde \(1\): mediante \((1,2)\) etiquetamos \(2\) con
    \((\color{#16a34a}{+1}, \, \color{#0284c7}{\min\{2, 1\} = 1})\).
  \item
    Desde \(2\): mediante \((2,4)\) etiquetamos \(4\) con
    \((\color{#16a34a}{+2}, \, \color{#0284c7}{\min\{1, 4 - 2\} = 1})\).
  \item
    Desde \(4\): etiquetar \(6\) con
    \((\color{#16a34a}{+4}, \, \color{#0284c7}{\min\{1, 5 - 3\} = 1})\).
  \item
    Desde \(6\): etiquetar sumidero \(t\) con
    \((\color{#16a34a}{+6}, \, \color{#0284c7}{\min\{1, 6 - 4\} = 1})\).
  \end{itemize}
\end{itemize}
\end{frame}

\begin{frame}{Ejemplo Etiquetado: Iteración 5 (Aumento y Actualización)}
\phantomsection\label{ejemplo-etiquetado-iteraciuxf3n-5-aumento-y-actualizaciuxf3n}
\begin{itemize}
\item
  \textbf{Camino de aumento detectado}:
  \[\mu_5 = (s \to 1 \to 2 \to 4 \to 6 \to t) \quad \text{con } \Delta_5 = \color{#0284c7}{\delta_t = 1}\]
\item
  \textbf{Actualización de caudales (\(\mathbf{f}^5\))}:

  \begin{itemize}
  \tightlist
  \item
    \(f_{s1} \leftarrow 5 + 1 = 6\).
  \item
    \(f_{12} \leftarrow 0 + 1 = 1 \implies\) Arco \((1,2)\)
    \textbf{saturado} (verde).
  \item
    \(f_{24} \leftarrow 2 + 1 = 3\), \(f_{46} \leftarrow 3 + 1 = 4\),
    \(f_{6t} \leftarrow 4 + 1 = 5\).
  \item
    Flujo de retorno circulante: \(f_0^5 = 7 + 1 = 8\).
  \end{itemize}
\end{itemize}
\end{frame}

\begin{frame}{Ejemplo Etiquetado: Iteración 6 (Criterio de Parada)}
\phantomsection\label{ejemplo-etiquetado-iteraciuxf3n-6-criterio-de-parada}
\begin{itemize}
\tightlist
\item
  \textbf{Propagación de marcas}:

  \begin{itemize}
  \tightlist
  \item
    Etiqueta en \(s\): \((\text{–}, \, \color{#0284c7}{+\infty})\).
  \item
    Desde \(s\): el arco \((s,2)\) está saturado (\(f_{s2} = 2\)).
    Etiquetar \(1\) con
    \((\color{#16a34a}{+s}, \, \color{#0284c7}{7 - 6 = 1})\).
  \item
    Desde \(1\): el arco \((1,3)\) está saturado (\(f_{13} = 5\)) y el
    arco \((1,2)\) está saturado (\(f_{12} = 1\)). No es posible
    etiquetar ningún nodo adicional desde \(1\).
  \end{itemize}
\item
  \textbf{Conclusión}: El sumidero \(t\) \textbf{no resulta etiquetado}
  (\(t \notin S\)) \(\implies\) \textbf{PARAR}.
\end{itemize}
\end{frame}

\begin{frame}{Ejemplo Etiquetado: Corte Mínimo Resultante}
\phantomsection\label{ejemplo-etiquetado-corte-muxednimo-resultante}
\begin{itemize}
\item
  \textbf{Solución Óptima}:

  \begin{itemize}
  \tightlist
  \item
    Valor del flujo máximo: \(f_0^* = 8\).
  \end{itemize}
\item
  \textbf{Subconjunto de vértices etiquetados}: \[N = S = \{s, 1\}\]
\item
  \textbf{Corte de capacidad mínima asociado}:
  \[w^+(N) = \{(s, 2), (1, 3), (1, 2)\}\]
\item
  \textbf{Capacidad del corte de mínima capacidad}:
  \[C(w^+(N)) = c_{s2} + c_{13} + c_{12} = 2 + 5 + 1 = 8 = f_0^*\]
\end{itemize}
\end{frame}

\begin{frame}{Ejemplo Etiquetado: Reproductor Interactivo}
\phantomsection\label{ejemplo-etiquetado-reproductor-interactivo}
\end{frame}

\begin{frame}{Ejemplo Etiquetado: Simbología y Código de Colores}
\phantomsection\label{ejemplo-etiquetado-simbologuxeda-y-cuxf3digo-de-colores}
\begin{itemize}
\tightlist
\item
  \textbf{Significado visual de la simbología en los gráficos}:

  \begin{itemize}
  \tightlist
  \item
    \textbf{Arcos en naranja (\(\color{#f97316}{\rightarrow}\))}: Camino
    de aumento seleccionado en la iteración actual.
  \item
    \textbf{Arcos en verde (\(\color{#16a34a}{\rightarrow}\))}: Arcos
    saturados a su capacidad máxima (\(f_u = c_u\)).
  \item
    \textbf{Arcos en rojo (\(\color{firebrick}{\rightarrow}\))}: Flujo
    parcial no nulo (\(0 < f_u < c_u\)).
  \item
    \textbf{Arcos en negro (\(\rightarrow\))}: Flujo nulo (\(f_u = 0\))
    o no transitado.
  \item
    \textbf{Nodos en naranja (\(\color{#f97316}{\bigcirc}\))}:
    Subconjunto de nodos etiquetados \(N = S\) que contienen a \(s\).
  \item
    \textbf{Nodos en verde (\(\color{#16a34a}{\bigcirc}\))}: Subconjunto
    de nodos no alcanzables \(V \setminus N\) que contienen a \(t\).
  \item
    \textbf{Arcos en azul (\(\color{#2563eb}{\rightarrow}\))}: Arcos
    salientes del \textbf{corte de capacidad mínima} \(w^+(N)\).
  \end{itemize}
\end{itemize}
\end{frame}

\begin{frame}{Teorema: Cota y Complejidad de Edmonds-Karp (1972)}
\phantomsection\label{teorema-cota-y-complejidad-de-edmonds-karp-1972}
\begin{tcolorbox}[enhanced jigsaw, coltitle=black, toprule=.15mm, opacitybacktitle=0.6, toptitle=1mm, leftrule=.75mm, breakable, left=2mm, opacityback=0, colback=white, bottomtitle=1mm, titlerule=0mm, title=\textcolor{quarto-callout-important-color}{\faExclamation}\hspace{0.5em}{Cota y complejidad de Edmonds-Karp y Dinitz (1970 / 1972)}, colbacktitle=quarto-callout-important-color!10!white, rightrule=.15mm, bottomrule=.15mm, arc=.35mm, colframe=quarto-callout-important-color-frame]

Si en el algoritmo de Ford-Fulkerson cada incremento de flujo se realiza
seleccionando sistemáticamente la cadena de aumento que posea el
\textbf{menor número de arcos} en el grafo residual \(G(\mathbf{f})\)
(garantizado mediante BFS), el flujo máximo se alcanza en no más de:

\[
\frac{1}{2} m n
\]

incrementos de flujo.

\end{tcolorbox}
\end{frame}

\begin{frame}{Importancia Teórica de Edmonds-Karp}
\phantomsection\label{importancia-teuxf3rica-de-edmonds-karp}
\begin{itemize}
\tightlist
\item
  \textbf{Garantía de Convergencia Polinomial}:

  \begin{itemize}
  \tightlist
  \item
    La propagación de marcas por niveles BFS en cada iteración requiere
    \(O(m)\) operaciones.
  \item
    La variante de Edmonds-Karp garantiza una complejidad temporal de
    \(O(m^2 n)\).
  \end{itemize}
\item
  \textbf{Independencia de Capacidades}:

  \begin{itemize}
  \tightlist
  \item
    Es totalmente independiente del valor numérico de las capacidades
    \(c_u\).
  \item
    Evita secuencias de iteraciones infinitas o no convergentes que
    pueden producirse en redes con capacidades irracionales si se eligen
    caminos arbitrarios (por ejemplo, mediante DFS).
  \end{itemize}
\end{itemize}
\end{frame}

\begin{frame}{Formulación Primal del Flujo Máximo}
\phantomsection\label{formulaciuxf3n-primal-del-flujo-muxe1ximo}
\begin{itemize}
\tightlist
\item
  \textbf{El problema del flujo máximo como programa lineal}:

  \begin{itemize}
  \tightlist
  \item
    Sea \((G = (V, U), \mathbf{c})\) una red de transporte con fuente
    \(s\) y sumidero \(t\).
  \item
    Variables: \(f_{ij} \ge 0\) para cada arco \((i,j) \in U\) y caudal
    de retorno \(f_0\).
  \end{itemize}
\item
  \textbf{Formulación del Modelo Primal}: \[
  \begin{aligned}
  \max_{f, f_0} \quad & f_0 \\
  \text{s.a.} \quad & \sum_{j \mid (j, i) \in U} f_{ji} - \sum_{j \mid (i, j) \in U} f_{ij} = \begin{cases} -f_0 & \text{si } i = s \\ 0 & \text{si } i \neq s, t \\ f_0 & \text{si } i = t \end{cases} \quad \forall i \in V \\
  & 0 \le f_{ij} \le c_{ij}, \quad \forall (i, j) \in U
  \end{aligned}
  \]
\end{itemize}
\end{frame}

\begin{frame}{Construcción del Modelo Dual}
\phantomsection\label{construcciuxf3n-del-modelo-dual}
\begin{itemize}
\tightlist
\item
  \textbf{Asignación de variables duales}:

  \begin{itemize}
  \tightlist
  \item
    Variable \(u_i \in \mathbb{R}\) por la conservación nodal en
    \(i \in V\) (potencial nodal).
  \item
    Variable \(w_{ij} \ge 0\) por la cota de capacidad
    \(f_{ij} \le c_{ij}\) en \((i, j) \in U\).
  \end{itemize}
\item
  \textbf{Objetivo dual}:

  \begin{itemize}
  \tightlist
  \item
    Minimizar el coste ponderado de las capacidades fijas de los arcos:
    \[\min_{u, w} \sum_{(i,j) \in U} c_{ij} w_{ij}\]
  \end{itemize}
\end{itemize}
\end{frame}

\begin{frame}{Fijación del Potencial de Referencia}
\phantomsection\label{fijaciuxf3n-del-potencial-de-referencia}
\begin{tcolorbox}[enhanced jigsaw, coltitle=black, toprule=.15mm, opacitybacktitle=0.6, toptitle=1mm, leftrule=.75mm, breakable, left=2mm, opacityback=0, colback=white, bottomtitle=1mm, titlerule=0mm, title=\textcolor{quarto-callout-note-color}{\faInfo}\hspace{0.5em}{Fijación de potencial de referencia}, colbacktitle=quarto-callout-note-color!10!white, rightrule=.15mm, bottomrule=.15mm, arc=.35mm, colframe=quarto-callout-note-color-frame]

Como la suma global de ecuaciones de conservación nodal es idénticamente
nula (\(\sum_{i \in V} b_i = 0\)), el sistema primal contiene 1 ecuación
redundante.

Fijando el potencial del sumidero en \(u_t = 0\), la condición de
dualidad sobre \(f_0\) impone: \[u_s - u_t = 1 \implies u_s = 1\]

\end{tcolorbox}
\end{frame}

\begin{frame}{Formulación Dual del Flujo Máximo}
\phantomsection\label{formulaciuxf3n-dual-del-flujo-muxe1ximo}
\begin{tcolorbox}[enhanced jigsaw, coltitle=black, toprule=.15mm, opacitybacktitle=0.6, toptitle=1mm, leftrule=.75mm, breakable, left=2mm, opacityback=0, colback=white, bottomtitle=1mm, titlerule=0mm, title=\textcolor{quarto-callout-note-color}{\faInfo}\hspace{0.5em}{Modelo lineal dual}, colbacktitle=quarto-callout-note-color!10!white, rightrule=.15mm, bottomrule=.15mm, arc=.35mm, colframe=quarto-callout-note-color-frame]

Teniendo en cuenta la normalización de potenciales \(u_s = 1, u_t = 0\),
el problema dual de programación lineal se formula como:

\[
\begin{aligned}
\min_{u, w} \quad & \sum_{(i,j) \in U} c_{ij} w_{ij} \\
\text{s.a.} \quad & u_j - u_i + w_{ij} \ge 0, \quad \forall (i, j) \in U \\
& u_s = 1, \quad u_t = 0 \\
& w_{ij} \ge 0, \quad \forall (i, j) \in U, \quad u_i \text{ libre}, \, \forall i \in V
\end{aligned}
\]

\end{tcolorbox}
\end{frame}

\begin{frame}{Interpretación Combinatoria del Modelo Dual}
\phantomsection\label{interpretaciuxf3n-combinatoria-del-modelo-dual}
\begin{itemize}
\tightlist
\item
  \textbf{Discretización binaria en vértices óptimos}:

  \begin{itemize}
  \tightlist
  \item
    En las soluciones básicas óptimas del dual, \(u_i \in \{0, 1\}\) y
    \(w_{ij} \in \{0, 1\}\).
  \end{itemize}
\item
  \textbf{Partición de los vértices (Corte)}:

  \begin{itemize}
  \tightlist
  \item
    \(u_i = 1\) define \(N = \{i \in V \mid u_i = 1\}\). Como
    \(u_s = 1\) y \(u_t = 0\), \((N, V \setminus N)\) es un
    \((s,t)\)-corte.
  \end{itemize}
\item
  \textbf{Identificación de arcos del corte}:

  \begin{itemize}
  \tightlist
  \item
    \(w_{ij} \ge u_i - u_j \implies w_{ij} = 1 \iff u_i = 1\) y
    \(u_j = 0\) (arcos salientes de \(N\)). Para los demás arcos,
    \(w_{ij} = 0\).
  \end{itemize}
\end{itemize}
\end{frame}

\begin{frame}{Interpretación de Dualidad Fuerte}
\phantomsection\label{interpretaciuxf3n-de-dualidad-fuerte}
\begin{tcolorbox}[enhanced jigsaw, coltitle=black, toprule=.15mm, opacitybacktitle=0.6, toptitle=1mm, leftrule=.75mm, breakable, left=2mm, opacityback=0, colback=white, bottomtitle=1mm, titlerule=0mm, title=\textcolor{quarto-callout-note-color}{\faInfo}\hspace{0.5em}{Interpretación combinatoria y dualidad fuerte}, colbacktitle=quarto-callout-note-color!10!white, rightrule=.15mm, bottomrule=.15mm, arc=.35mm, colframe=quarto-callout-note-color-frame]

Por el \textbf{Teorema de Dualidad Fuerte de la Programación Lineal}, el
valor óptimo del modelo primal \(\max f_0^*\) coincide exactamente con
el valor óptimo del modelo dual \(\min \sum c_{ij} w_{ij} = C(w^+(N))\).

Esto demuestra algebraicamente que \textbf{el Teorema de Flujo Máximo -
Corte Mínimo es la manifestación directa de la dualidad fuerte en
optimización lineal}.

\end{tcolorbox}
\end{frame}

\begin{frame}{Ejemplo Dualidad: Red de 3 Nodos}
\phantomsection\label{ejemplo-dualidad-red-de-3-nodos}
\begin{itemize}
\tightlist
\item
  Considérese la red de 3 nodos \(V = \{s, 1, t\}\) con arcos y
  capacidades:

  \begin{itemize}
  \tightlist
  \item
    \(c_{s1} = 3, \quad c_{st} = 2, \quad c_{1t} = 4\)
  \end{itemize}
\end{itemize}
\end{frame}

\begin{frame}{Ejemplo Dualidad: Planteamiento Primal y Dual}
\phantomsection\label{ejemplo-dualidad-planteamiento-primal-y-dual}
\begin{itemize}
\item
  \textbf{Modelo Primal}:
  \[\begin{aligned} \max_{x, f_0} \quad & f_0 \\ \text{s.a.} \quad & -x_{s1} - x_{st} = -f_0, \quad x_{s1} - x_{1t} = 0, \quad x_{st} + x_{1t} = f_0 \\ & 0 \le x_{s1} \le 3, \quad 0 \le x_{st} \le 2, \quad 0 \le x_{1t} \le 4 \end{aligned}\]
\item
  \textbf{Modelo Dual} (\(u_s = 1, u_t = 0\)):
  \[\begin{aligned} \min_{u, w} \quad & 3 w_{s1} + 2 w_{st} + 4 w_{1t} \\ \text{s.a.} \quad & u_1 + w_{s1} \ge 1, \quad w_{st} \ge 1, \quad -u_1 + w_{1t} \ge 0, \quad w_{ij} \ge 0 \end{aligned}\]
\end{itemize}
\end{frame}

\begin{frame}{Ejemplo Dualidad: Solución Óptima Primal y Dual}
\phantomsection\label{ejemplo-dualidad-soluciuxf3n-uxf3ptima-primal-y-dual}
\begin{itemize}
\tightlist
\item
  \textbf{Solución Primal Óptima}:

  \begin{itemize}
  \tightlist
  \item
    Flujo máximo \(f_0^* = 5\) con \(x_{s1}^* = 3\), \(x_{st}^* = 2\),
    \(x_{1t}^* = 3\).
  \end{itemize}
\item
  \textbf{Solución Dual Óptima}:

  \begin{itemize}
  \tightlist
  \item
    Potenciales \(u_s = 1, u_1 = 0, u_t = 0\).
  \item
    Multiplicadores \(w_{s1} = 1, w_{st} = 1, w_{1t} = 0\).
  \item
    Valor dual óptimo: \(3(1) + 2(1) + 4(0) = 5 = f_0^*\).
  \item
    Identifica el corte mínimo \(N = \{s\}\) con arcos salientes
    \(w^+(N) = \{(s,1), (s,t)\}\) de capacidad \(3 + 2 = 5\).
  \end{itemize}
\end{itemize}
\end{frame}

\begin{frame}{Ejemplo Dualidad: Representación de la Solución}
\phantomsection\label{ejemplo-dualidad-representaciuxf3n-de-la-soluciuxf3n}
\begin{itemize}
\tightlist
\item
  Visualización del flujo máximo de \(f_0^* = 5\) y la barrera del corte
  de capacidad mínima \(N = \{s\}\).
\end{itemize}
\end{frame}

\section{Problema de flujo a coste
mínimo}\label{problema-de-flujo-a-coste-muxednimo}

\begin{frame}{Introducción al Flujo a Coste Mínimo}
\phantomsection\label{introducciuxf3n-al-flujo-a-coste-muxednimo}
\begin{itemize}
\tightlist
\item
  \textbf{Más allá del flujo máximo}:

  \begin{itemize}
  \tightlist
  \item
    En logística y transporte no solo importa enviar la máxima cantidad,
    sino el \textbf{coste unitario de transporte}
    \(p_{ij} \in \mathbb{R}\) por cada unidad circulante.
  \end{itemize}
\item
  \textbf{Objetivo de optimización}:

  \begin{itemize}
  \tightlist
  \item
    Minimizar el coste económico total del envío manteniendo un volumen
    de transporte requerido \(\nu\) y respetando la capacidad física de
    los arcos:
  \end{itemize}
\end{itemize}

\[\min_{f} \sum_{(i,j) \in U} p_{ij} f_{ij}\]
\end{frame}

\begin{frame}{Planteamiento: Asignación de Tareas con Costes}
\phantomsection\label{planteamiento-asignaciuxf3n-de-tareas-con-costes}
\begin{itemize}
\tightlist
\item
  \textbf{Escenario industrial}:

  \begin{itemize}
  \tightlist
  \item
    \(n\) máquinas \(m_1 \dots m_n\) y \(n\) tareas \(t_1 \dots t_n\).
  \end{itemize}
\item
  \textbf{Limitación del enfoque de flujo máximo}:

  \begin{itemize}
  \tightlist
  \item
    Determina cuántas tareas pueden ejecutarse simultáneamente, pero
    ignora los distintos costes de procesamiento \(p_{ij}\) de cada
    máquina al realizar cada tarea.
  \end{itemize}
\item
  \textbf{Necesidad de la optimización a coste mínimo}:

  \begin{itemize}
  \tightlist
  \item
    Seleccionar la asignación completa que minimice el coste total
    invertido.
  \end{itemize}
\end{itemize}
\end{frame}

\begin{frame}{Asignación Bipartita en Red}
\phantomsection\label{asignaciuxf3n-bipartita-en-red}
\begin{itemize}
\tightlist
\item
  Representación en red bipartita con capacidades unitarias y costes de
  procesamiento \(p_{ij}\).
\end{itemize}
\end{frame}

\begin{frame}{Formulación del Flujo a Coste Mínimo}
\phantomsection\label{formulaciuxf3n-del-flujo-a-coste-muxednimo}
\begin{tcolorbox}[enhanced jigsaw, coltitle=black, toprule=.15mm, opacitybacktitle=0.6, toptitle=1mm, leftrule=.75mm, breakable, left=2mm, opacityback=0, colback=white, bottomtitle=1mm, titlerule=0mm, title=\textcolor{quarto-callout-note-color}{\faInfo}\hspace{0.5em}{Modelo PL del Flujo a Coste Mínimo}, colbacktitle=quarto-callout-note-color!10!white, rightrule=.15mm, bottomrule=.15mm, arc=.35mm, colframe=quarto-callout-note-color-frame]

Dada una red \((G = (V, U), \mathbf{c})\), costes unitarios
\(p_{ij} \in \mathbb{R}\) y un caudal total requerido \(f_0 = \nu\):

\[
\begin{aligned}
\min_{f} \quad & \sum_{(i,j) \in U} p_{ij} f_{ij} \\
\text{sujeto a} \quad & \sum_{j \mid (j, i) \in U} f_{ji} - \sum_{j \mid (i, j) \in U} f_{ij} = \begin{cases} -\nu & \text{si } i = s \\ 0 & \text{si } i \neq s, t \\ \nu & \text{si } i = t \end{cases} \quad \forall i \in V \\
& 0 \le f_{ij} \le c_{ij}, \quad \forall (i, j) \in U
\end{aligned}
\]

\end{tcolorbox}
\end{frame}

\begin{frame}{Ejemplo Coste Mínimo: Red de 5 Nodos}
\phantomsection\label{ejemplo-coste-muxednimo-red-de-5-nodos}
\begin{itemize}
\tightlist
\item
  Red de 5 nodos \(V = \{s, a, b, c, t\}\) con caudal requerido
  \(\nu = 4\).
\item
  Capacidades en azul \((c_{ij})\) y costes unitarios en rojo
  \([p_{ij}]\):
\end{itemize}
\end{frame}

\begin{frame}{Ejemplo Coste Mínimo: Comparación de Soluciones Factibles}
\phantomsection\label{ejemplo-coste-muxednimo-comparaciuxf3n-de-soluciones-factibles}
\begin{itemize}
\tightlist
\item
  Comparativa visual entre dos flujos compatibles de caudal 4:
  \(\mathbf{f}^A\) (coste 19) y \(\mathbf{f}^B\) (coste 18).
\end{itemize}
\end{frame}

\begin{frame}{Ejemplo Coste Mínimo: Evaluación de Flujos fA y fB}
\phantomsection\label{ejemplo-coste-muxednimo-evaluaciuxf3n-de-flujos-fa-y-fb}
\begin{itemize}
\tightlist
\item
  \textbf{Flujo Compatible \(\mathbf{f}^A\)} (en rojo):

  \begin{itemize}
  \tightlist
  \item
    Asignación:
    \(f_{sa}^A=1, f_{sb}^A=3, f_{ac}^A=1, f_{bc}^A=1, f_{bt}^A=2, f_{ct}^A=2\).
  \item
    Coste Total: \(c_A = 1(3) + 3(1) + 1(2) + 1(5) + 2(2) + 2(1) = 19\).
  \end{itemize}
\item
  \textbf{Flujo Compatible \(\mathbf{f}^B\)} (en azul):

  \begin{itemize}
  \tightlist
  \item
    Asignación:
    \(f_{sa}^B=2, f_{sb}^B=2, f_{ac}^B=2, f_{bc}^B=0, f_{bt}^B=2, f_{ct}^B=2\).
  \item
    Coste Total: \(c_B = 2(3) + 2(1) + 2(2) + 0(5) + 2(2) + 2(1) = 18\).
  \end{itemize}
\end{itemize}
\end{frame}

\begin{frame}{Ejemplo Coste Mínimo: Análisis del Ciclo de Diferencia}
\phantomsection\label{ejemplo-coste-muxednimo-anuxe1lisis-del-ciclo-de-diferencia}
\begin{itemize}
\tightlist
\item
  \textbf{Diferencia entre soluciones}:

  \begin{itemize}
  \tightlist
  \item
    \(\mathbf{\mu}^1 = \mathbf{f}^A - \mathbf{f}^B = (-1, 1, -1, 1, 0, 0) \implies\)
    Circuito cerrado \(s \to b \to c \to a \to s\).
  \end{itemize}
\item
  \textbf{Efecto de la reducción de coste}:

  \begin{itemize}
  \tightlist
  \item
    Sumar el ciclo inverso \(-\mathbf{\mu}^1\) reduce el coste en
    \(c_A - c_B = 1\).
  \end{itemize}
\item
  \textbf{Principio fundamental}:

  \begin{itemize}
  \tightlist
  \item
    La diferencia entre dos flujos cualesquiera de igual valor \(\nu\)
    es siempre una \textbf{combinación lineal de circuitos residuales}.
  \end{itemize}
\end{itemize}
\end{frame}

\begin{frame}{Ejemplo Coste Mínimo: Evaluación de Cadenas desde Flujo
Nulo}
\phantomsection\label{ejemplo-coste-muxednimo-evaluaciuxf3n-de-cadenas-desde-flujo-nulo}
\begin{itemize}
\tightlist
\item
  Evaluamos las 3 cadenas orientadas simples de \(s\) a \(t\) en la red
  nula \(\mathbf{f}^0 = \mathbf{0}\):

  \begin{enumerate}
  \tightlist
  \item
    \textbf{Cadena 1} (\(s \to a \to c \to t\)): Coste unitario
    \(3 + 2 + 1 = 6\).
  \item
    \textbf{Cadena 2} (\(s \to b \to t\)): Coste unitario \(1 + 2 = 3\).
  \item
    \textbf{Cadena 3} (\(s \to b \to c \to t\)): Coste unitario
    \(1 + 5 + 1 = 7\).
  \end{enumerate}
\end{itemize}
\end{frame}

\begin{frame}{Ejemplo Coste Mínimo: Solución Óptima Absoluta}
\phantomsection\label{ejemplo-coste-muxednimo-soluciuxf3n-uxf3ptima-absoluta}
\begin{itemize}
\tightlist
\item
  \textbf{Cadena de coste mínimo}:

  \begin{itemize}
  \tightlist
  \item
    \(s \to b \to t\) con coste unitario 3 y capacidad residual
    \(\min(4, 4) = 4\).
  \end{itemize}
\item
  \textbf{Asignación óptima}:

  \begin{itemize}
  \tightlist
  \item
    Canalizar las 4 unidades íntegramente por la trayectoria
    \(s \to b \to t\).
  \end{itemize}
\item
  \textbf{Coste mínimo absoluto}:

  \begin{itemize}
  \tightlist
  \item
    \(4 \times 3 = 12\) (estrictamente menor que los costes 19 y 18 de
    las soluciones de prueba \(\mathbf{f}^A\) y \(\mathbf{f}^B\)).
  \end{itemize}
\end{itemize}
\end{frame}

\begin{frame}{Coste y Capacidad Residual de Cadenas}
\phantomsection\label{coste-y-capacidad-residual-de-cadenas}
\begin{itemize}
\item
  \textbf{Coste Total Residual de una Cadena \(\mu\)}:
  \[p(\mu) = \sum_{u \in \mu^+} p_u - \sum_{u \in \mu^-} p_u\]
\item
  \textbf{Cadena Incremental}:

  \begin{itemize}
  \tightlist
  \item
    Cadena que admite un incremento positivo de flujo \(\delta > 0\) sin
    violar la no negatividad ni superar la capacidad física de ningún
    arco (\(f_u < c_u, \forall u \in \mu^+\) y
    \(f_u > 0, \forall u \in \mu^-\)).
  \end{itemize}
\end{itemize}
\end{frame}

\begin{frame}{Teorema: Optimalidad por Ausencia de Circuitos Negativos}
\phantomsection\label{teorema-optimalidad-por-ausencia-de-circuitos-negativos}
\begin{tcolorbox}[enhanced jigsaw, coltitle=black, toprule=.15mm, opacitybacktitle=0.6, toptitle=1mm, leftrule=.75mm, breakable, left=2mm, opacityback=0, colback=white, bottomtitle=1mm, titlerule=0mm, title=\textcolor{quarto-callout-important-color}{\faExclamation}\hspace{0.5em}{Optimalidad por ausencia de circuitos incrementales de coste negativo}, colbacktitle=quarto-callout-important-color!10!white, rightrule=.15mm, bottomrule=.15mm, arc=.35mm, colframe=quarto-callout-important-color-frame]

Un flujo compatible \(\mathbf{f}\) de valor \(f_0 = \nu\) es de
\textbf{coste mínimo} si y solo si la red residual \(G(\mathbf{f})\)
\textbf{no contiene ningún circuito incremental de coste total negativo}
(sin incluir el arco ficticio de retorno \((t, s)\)).

\end{tcolorbox}

\begin{itemize}
\tightlist
\item
  \textbf{Interpretación}: Si existiera en \(G(\mathbf{f})\) un circuito
  con \(p(\mu) < 0\), aumentar flujo a su través mantendría el caudal
  \(\nu\) pero reduciría estrictamente el coste total, contradiciendo la
  optimalidad.
\end{itemize}
\end{frame}

\begin{frame}{Teorema: Incrementos por Caminos de Coste Mínimo}
\phantomsection\label{teorema-incrementos-por-caminos-de-coste-muxednimo}
\begin{tcolorbox}[enhanced jigsaw, coltitle=black, toprule=.15mm, opacitybacktitle=0.6, toptitle=1mm, leftrule=.75mm, breakable, left=2mm, opacityback=0, colback=white, bottomtitle=1mm, titlerule=0mm, title=\textcolor{quarto-callout-important-color}{\faExclamation}\hspace{0.5em}{Incrementos por caminos de coste mínimo}, colbacktitle=quarto-callout-important-color!10!white, rightrule=.15mm, bottomrule=.15mm, arc=.35mm, colframe=quarto-callout-important-color-frame]

Dado un flujo \(\mathbf{f}\) de valor \(\nu\) que es de coste mínimo, si
se incrementa el flujo en \(\delta\) unidades a lo largo de la cadena
incremental de \(s\) a \(t\) que posea el \textbf{coste total residual
mínimo} en \(G(\mathbf{f})\), el flujo resultante es de coste mínimo
para el nuevo caudal \(\nu + \delta\).

\end{tcolorbox}

\begin{itemize}
\tightlist
\item
  \textbf{Garantía de construcción}: Valida la estrategia de avanzar
  iterativamente incrementando el caudal a través de las trayectorias
  residuales de menor coste.
\end{itemize}
\end{frame}

\begin{frame}{Estrategia Algorítmica de Resolución Iterativa}
\phantomsection\label{estrategia-algoruxedtmica-de-resoluciuxf3n-iterativa}
\begin{tcolorbox}[enhanced jigsaw, coltitle=black, toprule=.15mm, opacitybacktitle=0.6, toptitle=1mm, leftrule=.75mm, breakable, left=2mm, opacityback=0, colback=white, bottomtitle=1mm, titlerule=0mm, title=\textcolor{quarto-callout-note-color}{\faInfo}\hspace{0.5em}{Estrategia algorítmica de resolución iterativa}, colbacktitle=quarto-callout-note-color!10!white, rightrule=.15mm, bottomrule=.15mm, arc=.35mm, colframe=quarto-callout-note-color-frame]

\begin{enumerate}
\tightlist
\item
  \textbf{Determinación del Flujo Inicial}:

  \begin{itemize}
  \tightlist
  \item
    Comenzar con \(\mathbf{f}^0 = \mathbf{0}\) para \(\nu = 0\). Si
    \(p_{ij} \ge 0\), no hay circuitos de coste negativo
    \(\implies \mathbf{f}^0\) es trivialmente óptimo.
  \end{itemize}
\item
  \textbf{Crecimiento Iterativo}:

  \begin{itemize}
  \tightlist
  \item
    Incrementar progresivamente el caudal \(\nu \to \nu + \delta\) a lo
    largo de la cadena de coste residual mínimo en \(G(\mathbf{f})\)
    hasta alcanzar el caudal objetivo \(\nu\).
  \end{itemize}
\end{enumerate}

\end{tcolorbox}
\end{frame}

\begin{frame}{Correspondencia de Circuitos Negativos en G(f)}
\phantomsection\label{correspondencia-de-circuitos-negativos-en-gf}
\begin{itemize}
\tightlist
\item
  \textbf{Red residual \(G(\mathbf{f})\) como grafo de distancias}:

  \begin{itemize}
  \tightlist
  \item
    Cada arco en \(G(\mathbf{f})\) lleva asignado el coste residual
    \(p_u\) (si es directo) o \(-p_u\) (si es inverso).
  \end{itemize}
\item
  \textbf{Equivalencia de circuitos}:

  \begin{itemize}
  \tightlist
  \item
    Un circuito de longitud total estrictamente negativa en
    \(G(\mathbf{f})\) se corresponde de forma unívoca con un
    \textbf{ciclo incremental de coste negativo} en la red de transporte
    original.
  \end{itemize}
\end{itemize}
\end{frame}

\begin{frame}{Esquema: Algoritmo de Flujo a Coste Mínimo (Pasos 1 y 2)}
\phantomsection\label{esquema-algoritmo-de-flujo-a-coste-muxednimo-pasos-1-y-2}
\begin{itemize}
\tightlist
\item
  \textbf{Paso 1 (Inicialización)}:

  \begin{itemize}
  \tightlist
  \item
    Obtener un flujo compatible inicial \(\mathbf{f}\) (típicamente
    \(\mathbf{f}^0 = \mathbf{0}\) para \(\nu = 0\)).
  \end{itemize}
\item
  \textbf{Paso 2 (Eliminación de circuitos de coste negativo)}:

  \begin{enumerate}
  \tightlist
  \item
    Construir el grafo de distancias residual \(G(\mathbf{f})\).
  \item
    Si \(G(\mathbf{f})\) contiene circuitos de longitud negativa,
    aumentar el flujo en su capacidad residual y repetir Paso 2.1.
  \item
    Si no existen circuitos negativos, ir al \textbf{Paso 3}.
  \end{enumerate}
\end{itemize}
\end{frame}

\begin{frame}{Esquema: Algoritmo de Flujo a Coste Mínimo (Paso 3)}
\phantomsection\label{esquema-algoritmo-de-flujo-a-coste-muxednimo-paso-3}
\begin{itemize}
\tightlist
\item
  \textbf{Paso 3 (Aumento del valor del flujo)}:

  \begin{itemize}
  \tightlist
  \item
    Si \(f_0 = \nu \implies\) \textbf{PARAR}: Solución óptima alcanzada.
  \item
    Si \(f_0 < \nu\), buscar en \(G(\mathbf{f})\) el camino de longitud
    mínima de \(s\) a \(t\):

    \begin{enumerate}
    \tightlist
    \item
      Si no existe camino \(\implies\) \textbf{PARAR}: Incompatible para
      el caudal \(\nu\).
    \item
      Si existe camino, calcular la capacidad residual \(c(s, t)\).
    \item
      Incrementar el flujo en \(\Delta = \min\{c(s, t), \nu - f_0\}\).
    \item
      Actualizar \(f_0 \leftarrow f_0 + \Delta\), reconstruir
      \(G(\mathbf{f})\) e ir al \textbf{Paso 3}.
    \end{enumerate}
  \end{itemize}
\end{itemize}
\end{frame}

\begin{frame}{Ejemplo Traza Coste Mínimo: Planteamiento de la Red}
\phantomsection\label{ejemplo-traza-coste-muxednimo-planteamiento-de-la-red}
\begin{itemize}
\tightlist
\item
  Red de 5 nodos \(V = \{s, a, b, c, t\}\) con 7 arcos dirigidos
  \((c_{ij})[p_{ij}]\):

  \begin{itemize}
  \tightlist
  \item
    \((s,a): (3)[3]\), \((s,b): (2)[3]\), \((a,c): (2)[-1]\),
    \((a,t): (4)[-1]\)
  \item
    \((b,a): (2)[-2]\), \((c,b): (2)[1]\), \((c,t): (5)[2]\)
  \end{itemize}
\item
  \textbf{Objetivo de optimización}:

  \begin{itemize}
  \tightlist
  \item
    Transmitir un caudal objetivo de \(\nu = 5\) unidades desde \(s\)
    hasta \(t\) al menor coste total posible.
  \end{itemize}
\end{itemize}
\end{frame}

\begin{frame}{Ejemplo Traza: Paso 1 (Inicialización con Flujo Nulo)}
\phantomsection\label{ejemplo-traza-paso-1-inicializaciuxf3n-con-flujo-nulo}
\begin{itemize}
\tightlist
\item
  \textbf{Solución inicial de flujo nulo}:

  \begin{itemize}
  \tightlist
  \item
    \(\mathbf{f}^0 = \mathbf{0}\) para caudal circulante \(f_0 = 0\).
  \item
    Coste acumulado inicial:
    \(C(\mathbf{f}^0) = \sum_{(i,j) \in U} p_{ij} f_{ij}^0 = 0\).
  \end{itemize}
\item
  \textbf{Justificación}:

  \begin{itemize}
  \tightlist
  \item
    El flujo nulo satisface trivialmente la conservación en todos los
    vértices (\(0 = 0\)) y la compatibilidad con las capacidades
    (\(0 \le c_{ij}\)).
  \end{itemize}
\end{itemize}
\end{frame}

\begin{frame}{Ejemplo Traza: Paso 2.1 y 2.2 (Detección del Circuito
Negativo)}
\phantomsection\label{ejemplo-traza-paso-2.1-y-2.2-detecciuxf3n-del-circuito-negativo}
\begin{itemize}
\tightlist
\item
  \textbf{Construcción del Grafo Residual \(G(\mathbf{f}^0)\)}:

  \begin{itemize}
  \tightlist
  \item
    Al ser \(\mathbf{f}^0 = \mathbf{0}\), todos los arcos admiten
    \(c_{ij} > 0\). Los pesos unitarios en \(G(\mathbf{f}^0)\) son:
    \[p_{sa}=3, \, p_{sb}=3, \, p_{ba}=-2, \, p_{ac}=-1, \, p_{cb}=1, \, p_{at}=-1, \, p_{ct}=2\]
  \end{itemize}
\item
  \textbf{Detección del Circuito Residual
  \(C_1 = a \to c \to b \to a\)}:

  \begin{itemize}
  \tightlist
  \item
    Coste total del circuito:
    \[p(C_1) = p_{ac} + p_{cb} + p_{ba} = (-1) + 1 + (-2) = -2 < 0\]
  \item
    Capacidad residual disponible en \(C_1\):
    \[\delta_1 = \min \{ c_{ac} - 0, \, c_{cb} - 0, \, c_{ba} - 0 \} = \min \{2, 2, 2\} = 2 \text{ unidades}\]
  \end{itemize}
\end{itemize}
\end{frame}

\begin{frame}{Ejemplo Traza: Paso 2.3 y 2.4 (Eliminación del Circuito
Negativo)}
\phantomsection\label{ejemplo-traza-paso-2.3-y-2.4-eliminaciuxf3n-del-circuito-negativo}
\begin{itemize}
\tightlist
\item
  \textbf{Incremento de flujo a lo largo de \(C_1\)}:

  \begin{itemize}
  \tightlist
  \item
    Asignamos 2 unidades en
    \(C_1 \implies f_{ac}^1 = 2, f_{cb}^1 = 2, f_{ba}^1 = 2\).
  \item
    Recálculo del coste acumulado:
    \[C(\mathbf{f}^1) = 2(-1) + 2(1) + 2(-2) = -4\]
  \item
    El coste baja de \(0\) a \(-4\) manteniendo el caudal entrante
    \(f_0 = 0\).
  \end{itemize}
\item
  \textbf{Nuevo Grafo Residual \(G(\mathbf{f}^1)\)}:

  \begin{itemize}
  \tightlist
  \item
    Los arcos \((a,c), (c,b), (b,a)\) se saturan (\(f=c=2\)),
    invirtiéndose a: \((c,a)\) con peso \(+1\), \((b,c)\) con peso
    \(-1\), y \((a,b)\) con peso \(+2\).
  \item
    ¡Ya no existen circuitos de coste negativo! \(\mathbf{f}^1\) es el
    punto de partida óptimo para \(f_0 = 0\).
  \end{itemize}
\end{itemize}
\end{frame}

\begin{frame}{Ejemplo Traza: Paso 3 - Iteración 1 (Camino Mínimo s
-\textgreater{} a -\textgreater{} t)}
\phantomsection\label{ejemplo-traza-paso-3---iteraciuxf3n-1-camino-muxednimo-s---a---t}
\begin{itemize}
\tightlist
\item
  \textbf{Caudal actual \(f_0 = 0 < 5 = \nu\)}:

  \begin{itemize}
  \tightlist
  \item
    Buscamos en \(G(\mathbf{f}^1)\) el camino de longitud residual
    mínima de \(s\) a \(t\).
  \end{itemize}
\item
  \textbf{Detección del Camino Mínimo \(\mu_1 = s \to a \to t\)}:

  \begin{itemize}
  \tightlist
  \item
    Coste residual total de la trayectoria:
    \[p(\mu_1) = p_{sa} + p_{at} = 3 + (-1) = 2\]
  \end{itemize}
\item
  \textbf{Cálculo de la capacidad residual e incremento}:

  \begin{itemize}
  \tightlist
  \item
    Capacidad residual de la cadena:
    \(c(\mu_1) = \min \{c_{sa} - 0, \, c_{at} - 0\} = \min \{3, 4\} = 3\)
    unidades.
  \item
    Incremento de flujo: \(\Delta_1 = \min \{3, \, 5 - 0\} = 3\)
    unidades.
  \end{itemize}
\end{itemize}
\end{frame}

\begin{frame}{Ejemplo Traza: Paso 3 - Iteración 1 (Actualización f0 =
3)}
\phantomsection\label{ejemplo-traza-paso-3---iteraciuxf3n-1-actualizaciuxf3n-f0-3}
\begin{itemize}
\tightlist
\item
  \textbf{Actualización del vector de flujo \(\mathbf{f}^2\)}:

  \begin{itemize}
  \tightlist
  \item
    Asignamos \(\Delta_1 = 3\) unidades en \(\mu_1\): \(f_{sa}^2 = 3\)
    (saturado), \(f_{at}^2 = 3\).
  \item
    Caudal circulante alcanzado: \(f_0 = 3\).
  \item
    Coste total acumulado: \(C(\mathbf{f}^2) = -4 + 3(2) = 2\).
  \end{itemize}
\item
  \textbf{Construcción del nuevo Grafo Residual \(G(\mathbf{f}^2)\)}:

  \begin{itemize}
  \tightlist
  \item
    Arco \((s,a)\) saturado (\(f_{sa}=3=c_{sa}\)) \(\implies\) genera
    arco inverso \((a,s)\) con peso \(-3\).
  \item
    Arco \((a,t)\) parcialmente cargado (\(f_{at}=3 < c_{at}=4\))
    \(\implies\) mantiene directo \((a,t)\) con peso \(-1\) e inverso
    \((t,a)\) con peso \(+1\).
  \end{itemize}
\end{itemize}
\end{frame}

\begin{frame}{Ejemplo Traza: Paso 3 - Iteración 2 (Camino Mínimo s
-\textgreater{} b -\textgreater{} c -\textgreater{} a -\textgreater{}
t)}
\phantomsection\label{ejemplo-traza-paso-3---iteraciuxf3n-2-camino-muxednimo-s---b---c---a---t}
\begin{itemize}
\tightlist
\item
  \textbf{Caudal actual \(f_0 = 3 < 5 = \nu\)}:

  \begin{itemize}
  \tightlist
  \item
    Buscamos en \(G(\mathbf{f}^2)\) el camino de longitud residual
    mínima de \(s\) a \(t\).
  \end{itemize}
\item
  \textbf{Detección del Camino Mínimo
  \(\mu_2 = s \to b \to c \to a \to t\)}:

  \begin{itemize}
  \tightlist
  \item
    La trayectoria utiliza los arcos directos \((s,b)\) y \((a,t)\), y
    los arcos inversos \((c,b)\) y \((a,c)\).
  \item
    Coste residual total de la trayectoria:
    \[p(\mu_2) = p_{sb} + (-p_{cb}) + (-p_{ac}) + p_{at} = 3 + (-1) + 1 + (-1) = 2 \quad (\text{o } 4 \text{ sin cancelación})\]
  \end{itemize}
\item
  \textbf{Cálculo de la capacidad residual e incremento}:

  \begin{itemize}
  \tightlist
  \item
    \(c(\mu_2) = \min \{ c_{sb}-0, \, f_{cb}^2, \, f_{ac}^2, \, c_{at}-3 \} = \min \{2, 2, 2, 1\} = 1\)
    unidad.
  \item
    Incremento de flujo: \(\Delta_2 = \min \{1, \, 5 - 3\} = 1\) unidad.
  \end{itemize}
\end{itemize}
\end{frame}

\begin{frame}{Ejemplo Traza: Paso 3 - Iteración 2 (Actualización f0 =
4)}
\phantomsection\label{ejemplo-traza-paso-3---iteraciuxf3n-2-actualizaciuxf3n-f0-4}
\begin{itemize}
\tightlist
\item
  \textbf{Actualización del vector de flujo \(\mathbf{f}^3\)}:

  \begin{itemize}
  \tightlist
  \item
    Flujo directo aumenta en 1 unidad: \(f_{sb}^3 = 1\),
    \(f_{at}^3 = 4\) (saturado).
  \item
    Flujo en arcos inversos se reduce en 1 unidad: \(f_{cb}^3 = 1\),
    \(f_{ac}^3 = 1\).
  \item
    Caudal circulante alcanzado: \(f_0 = 4\).
  \item
    Coste total acumulado: \(C(\mathbf{f}^3) = 2 + 1(2) = 4\).
  \end{itemize}
\item
  \textbf{Construcción del nuevo Grafo Residual \(G(\mathbf{f}^3)\)}:

  \begin{itemize}
  \tightlist
  \item
    Arco \((a,t)\) queda saturado (\(f_{at}=4=c_{at}\)) \(\implies\)
    desaparece su arco directo y permanece únicamente el inverso
    \((t,a)\) con peso \(+1\).
  \end{itemize}
\end{itemize}
\end{frame}

\begin{frame}{Ejemplo Traza: Paso 3 - Iteración 3 (Camino Mínimo s
-\textgreater{} b -\textgreater{} c -\textgreater{} t)}
\phantomsection\label{ejemplo-traza-paso-3---iteraciuxf3n-3-camino-muxednimo-s---b---c---t}
\begin{itemize}
\tightlist
\item
  \textbf{Caudal actual \(f_0 = 4 < 5 = \nu\)}:

  \begin{itemize}
  \tightlist
  \item
    Buscamos en \(G(\mathbf{f}^3)\) el camino de longitud residual
    mínima de \(s\) a \(t\).
  \end{itemize}
\item
  \textbf{Detección del Camino Mínimo \(\mu_3 = s \to b \to c \to t\)}:

  \begin{itemize}
  \tightlist
  \item
    La trayectoria utiliza los arcos directos \((s,b)\) y \((c,t)\), y
    el arco inverso \((c,b)\).
  \item
    Coste residual total de la trayectoria:
    \[p(\mu_3) = p_{sb} + (-p_{cb}) + p_{ct} = 3 + (-1) + 2 = 4\]
  \end{itemize}
\item
  \textbf{Cálculo de la capacidad residual e incremento final}:

  \begin{itemize}
  \tightlist
  \item
    \(c(\mu_3) = \min \{ c_{sb}-1, \, f_{cb}^3, \, c_{ct}-0 \} = \min \{1, 1, 5\} = 1\)
    unidad.
  \item
    Incremento final: \(\Delta_3 = \min \{1, \, 5 - 4\} = 1\) unidad.
  \end{itemize}
\end{itemize}
\end{frame}

\begin{frame}{Ejemplo Traza: Paso 3 - Iteración 3 (Actualización f0 = 5
y Parada)}
\phantomsection\label{ejemplo-traza-paso-3---iteraciuxf3n-3-actualizaciuxf3n-f0-5-y-parada}
\begin{itemize}
\tightlist
\item
  \textbf{Actualización final de flujos}:

  \begin{itemize}
  \tightlist
  \item
    \(f_{sb}^* = 1 + 1 = 2\) (saturando el arco \((s,b)\)).
  \item
    \(f_{cb}^* = 1 - 1 = 0\) (se deshace por completo el flujo residual
    en \((c,b)\)).
  \item
    \(f_{ct}^* = 0 + 1 = 1\).
  \end{itemize}
\item
  \textbf{Condición de Parada Cumplida}:

  \begin{itemize}
  \tightlist
  \item
    El caudal circulante alcanza la meta objetivo \(f_0 = 5 = \nu\).
  \item
    \textbf{PARAR}: Solución óptima global alcanzada.
  \end{itemize}
\end{itemize}
\end{frame}

\begin{frame}{Ejemplo Traza: Solución Óptima Final y Coste Mínimo Total}
\phantomsection\label{ejemplo-traza-soluciuxf3n-uxf3ptima-final-y-coste-muxednimo-total}
\begin{itemize}
\item
  \textbf{Distribución de flujo óptima final para \(\nu = 5\) unidades}:
  \[f_{sa}^* = 3, \quad f_{sb}^* = 2, \quad f_{ba}^* = 2, \quad f_{ac}^* = 1, \quad f_{cb}^* = 0, \quad f_{at}^* = 4, \quad f_{ct}^* = 1\]
\item
  \textbf{Cálculo del Coste Mínimo Óptimo Total (\(C^*\))}: \[
  \begin{aligned}
  C^* &= 3(p_{sa}) + 2(p_{sb}) + 2(p_{ba}) + 1(p_{ac}) + 0(p_{cb}) + 4(p_{at}) + 1(p_{ct}) \\
  &= 3(3) + 2(3) + 2(-2) + 1(-1) + 0(1) + 4(-1) + 1(2) \\
  &= 9 + 6 - 4 - 1 + 0 - 4 + 2 = 8
  \end{aligned}
  \]
\end{itemize}
\end{frame}

\begin{frame}{Ejemplo Traza: Reproductor Interactivo}
\phantomsection\label{ejemplo-traza-reproductor-interactivo}
\end{frame}

\section{El problema del transbordo}\label{el-problema-del-transbordo}

\begin{frame}{El Problema del Transbordo: Introducción}
\phantomsection\label{el-problema-del-transbordo-introducciuxf3n}
\begin{itemize}
\tightlist
\item
  \textbf{Extensión a redes complejas}:

  \begin{itemize}
  \tightlist
  \item
    Generaliza el flujo de coste mínimo de un par \((s, t)\) a redes con
    \textbf{múltiples nodos de oferta y demanda}.
  \item
    Formulación lineal más general e integradora dentro de la
    optimización en redes.
  \end{itemize}
\item
  \textbf{Clasificación de balances nodales (\(b_i\))}:

  \begin{itemize}
  \tightlist
  \item
    \textbf{Oferta (\(b_i > 0\))}: El nodo genera o dispone de \(b_i\)
    unidades de mercancía.
  \item
    \textbf{Demanda (\(b_i < 0\))}: El nodo requiere o consume \(|b_i|\)
    unidades de mercancía.
  \item
    \textbf{Transbordo (\(b_i = 0\))}: El nodo no produce ni consume;
    punto intermedio de paso (\(\sum f_{in} = \sum f_{out}\)).
  \end{itemize}
\item
  \textbf{Hipótesis de Equilibrio Global}: \(\sum_{i=1}^n b_i = 0\)
\end{itemize}
\end{frame}

\begin{frame}{Modelo Matemático del Transbordo}
\phantomsection\label{modelo-matemuxe1tico-del-transbordo}
\begin{tcolorbox}[enhanced jigsaw, coltitle=black, toprule=.15mm, opacitybacktitle=0.6, toptitle=1mm, leftrule=.75mm, breakable, left=2mm, opacityback=0, colback=white, bottomtitle=1mm, titlerule=0mm, title=\textcolor{quarto-callout-note-color}{\faInfo}\hspace{0.5em}{Modelo lineal del problema del transbordo}, colbacktitle=quarto-callout-note-color!10!white, rightrule=.15mm, bottomrule=.15mm, arc=.35mm, colframe=quarto-callout-note-color-frame]

Dada una red dirigida \(G = (V, U)\) con capacidades
\([\ell_{ij}, c_{ij}]\) y costes unitarios \(p_{ij}\):

\[
\begin{aligned}
\min_{f \ge 0} \quad & \sum_{(i,j) \in U} p_{ij} f_{ij} \\
\text{sujeto a} \quad & \sum_{j \mid (i, j) \in U} f_{ij} - \sum_{k \mid (k, i) \in U} f_{ki} = b_i, \quad \forall i \in V \\
& \ell_{ij} \le f_{ij} \le c_{ij}, \quad \forall (i, j) \in U
\end{aligned}
\]

\end{tcolorbox}
\end{frame}

\begin{frame}{El Transbordo como Marco Unificador}
\phantomsection\label{el-transbordo-como-marco-unificador}
\begin{enumerate}
\tightlist
\item
  \textbf{Problema de Transporte}: Nodos divididos estrictamente en
  orígenes (\(b_i > 0\)) y destinos (\(b_i < 0\)), sin nodos de
  transbordo.
\item
  \textbf{Problema de Asignación}: Caso de transporte bipartito con
  \(b_i = +1\), \(b_j = -1\) y capacidades unitarias.
\item
  \textbf{Problema del Camino Mínimo}: Red con un único origen
  \(b_s = +1\), un único destino \(b_t = -1\) y transbordos \(b_i = 0\).
\item
  \textbf{Problema del Flujo Máximo}: Red con arco de retorno ficticio
  \((t, s)\) e igualación de balances a cero (\(b_i = 0\)).
\item
  \textbf{Problema del Flujo a Coste Mínimo}: Formulación del transbordo
  con fuente \(b_s = \nu\), sumidero \(b_t = -\nu\) y transbordos
  \(b_i = 0\).
\end{enumerate}
\end{frame}




\end{document}
